\documentclass[11pt,a4paper]{article}
\usepackage[margin=27mm,headheight=14pt]{geometry}
\usepackage[T1]{fontenc}
\usepackage{lmodern,amsmath,amssymb,amsthm,mathtools,mathrsfs}
\usepackage{microtype,booktabs,longtable,array,enumitem,xcolor,graphicx}
\usepackage{xurl}
\usepackage[colorlinks=true,allcolors=blue!45!black]{hyperref}
\usepackage{bookmark,needspace}
\setlength{\emergencystretch}{2em}
\setlength{\parskip}{3pt}
\setlist{itemsep=2pt,topsep=4pt}
\allowdisplaybreaks[2]
\newtheorem{theorem}{Theorem}[section]
\newtheorem{proposition}[theorem]{Proposition}
\newtheorem{lemma}[theorem]{Lemma}
\newtheorem{corollary}[theorem]{Corollary}
\newtheorem{counterexample}[theorem]{Counterexample}
\theoremstyle{definition}
\newtheorem{definition}[theorem]{Definition}
\newtheorem{remark}[theorem]{Remark}
\newcommand{\Sh}{\operatorname{Sh}}
\newcommand{\Hom}{\operatorname{Hom}}
\newcommand{\Map}{\operatorname{Map}}
\newcommand{\Pro}{\operatorname{Pro}}
\newcommand{\colim}{\operatorname*{colim}}
\newcommand{\limm}{\operatorname*{lim}}
\newcommand{\id}{\operatorname{id}}
\newcommand{\Spec}{\operatorname{Spec}}
\newcommand{\im}{\operatorname{im}}
\newcommand{\cl}{\operatorname{cl}}
\newcommand{\R}{\mathbb R}
\newcommand{\C}{\mathbb C}
\newcommand{\Z}{\mathbb Z}
\newcommand{\Q}{\mathbb Q}
\newcommand{\F}{\mathbb F}
\newcommand{\recordurl}{https://github.com/KokunoYumeto/mathematics-workbenches}
\newcommand{\workbenchurl}{https://github.com/KokunoYumeto/homotopy-and-categories-workbench}
\title{Grothendieck's tensor divided-power categories\\ and cellwise dessins after Mulase and Penkava}
\author{}
\date{23 November 2026}
\begin{document}
\maketitle
\begin{abstract}
Grothendieck's cote~137 proposes tensor divided-power data as a description of polynomial categories and asks for additional structure recovering cochains. We separate homogeneous substitution from the mixed-degree coherence needed for polynomial composition, give explicit counterexamples to the bare equivalence, and formulate a finite coherence repair. For the walking-complex construction over a commutative ring, whose equivalence for targets with kernels is known, we isolate the precise matrix kernels required in an arbitrary additive target. We recover the full cosimplicial ring of integral cochains from its natural divided-power action on finite test objects. We characterize bounded normalization by uniformly bounded polynomial fibre ranks and distinguish it from derived perfection. For clean dessins, we extend Mulase and Penkava's classification of edge-fixed ribbon symmetries to the full holomorphic group preserving every ordinary cell, including the low-valence sphere cases. We distinguish labelled-cell counts from rigidity and prove rigidity for noncompact cellular maps. Exact Kummer images describe the associated diagonal Fermat fibres. The mathematics and text were produced by GPT-6 Astra (ChatGPT Pro mode), whose checks are individually recorded; no human has checked the work.
\end{abstract}

\section{Polynomial information and the questions in the manuscripts}
In cote~137, archive sheet~38 (PDF page~39), Grothendieck asks how the
additional structure on cochains ``permet de reconstituer le type
d'homotopie du $X_*$''. The preceding sheet calls it ``structure
suppl\'ementaire''. We make the reading precise: it is the entire
natural divided-power action on finite test objects. His later
axioms for tensor divided-power categories are on archive
sheets~80--85 (PDF pages~81--86). They raise a distinct question,
whether homogeneous substitution already determines composition
of arbitrary polynomial maps \cite{Groth}.

This section explains the distinctions needed to state the results. A homogeneous polynomial law of degree $d$ from a $k$-module $V$ to a $k$-module $W$ is a family of maps $R\otimes_kV\to R\otimes_kW$, natural in commutative $k$-algebras $R$, and homogeneous of degree $d$. Its universal representing module is $\Gamma^dV$; this is Roby's divided-power construction \cite[Theorem~IV.1, p.~266]{Roby}. Naturality after scalar extension is part of the definition. Agreement at the points of a finite field alone is weaker.

For finite projective modules the category used for strict polynomial functors has morphisms $\Gamma^d\Hom_k(V,W)$, with a fixed degree $d$. Its linear representations are the strict polynomial functors of degree $d$ \cite[Section~2, pp.~3--5]{Krause}. This category differs from the category with degree-$d$ polynomial laws $V\rightsquigarrow W$ as morphisms. The latter has morphisms $\Hom_k(\Gamma^dV,W)$, and degrees multiply under composition. The issue in cote~137, PDF pages~81--86, is how the homogeneous pieces fit together into this second, polynomial category.

Call \emph{bare divided-power data} an additive $k$-linear symmetric monoidal category $\mathcal C$, a coherent exponential functor $G^*$ with $G^0=\mathbf1$, $G^1=\id$, the scalar rule $G^n(\lambda\id_X)=\lambda^n\id_{G^nX}$, and natural maps
\[
 \alpha_{p,q}:G^{pq}\longrightarrow G^pG^q
\]
satisfying
\begin{equation}\label{eq:homogeneous}
 G^p(\alpha_{q,r})\alpha_{p,qr}
 =\alpha_{p,q}^{G^r}\alpha_{pq,r}.
\end{equation}
The one-index maps are identities and $\alpha_{0,q}$ is the identity of the unit. We state explicitly when the additional maps $\alpha_{q,0}$ are required to be invertible.

\begin{theorem}[The two missing requirements]\label{thm:bare}
Invertibility of $\alpha_{q,0}$ is independent of the other bare axioms. Even when all these maps are invertible, the bare axioms do not force polynomial composition compatible with the specified exponential and tensor-unit comparison maps. Both failures occur in finite-dimensional vector spaces, over $\F_2$ and $\F_3$ respectively.
\end{theorem}

The first counterexample concerns an omitted unit condition; the second concerns compatibility between different degrees. Their distinction is essential: adding the first condition does not repair the second failure. The next theorem gives a precise composition criterion using only finite sums, natural maps and equations; Section~3 proves it. Its associativity calculation is the finite-bound form of the classical co-Kleisli construction, dual to \cite[Definition~5.2.10, p.~194]{Riehl}; the extra point is the prescribed mixed-degree normalization.

The precise coherence criterion uses only finite divided-power sums. Put $B_mX=\bigoplus_{i=0}^mG^iX$. For $d,m\ge0$, the exponential decomposition prescribes a map
\[
 D_{d,m}:B_{dm}X\longrightarrow B_d(B_mX).
\]
In the summand of $G^q(B_mX)$ indexed by $(\nu_0,\ldots,\nu_m)$ with $\sum\nu_i=q\le d$, its weight is $n=\sum i\nu_i$. It is zero on other source degrees; on $G^nX$ it is
\[
 G^nX\xrightarrow{\Delta}\bigotimes_iG^{i\nu_i}X
 \xrightarrow{\otimes_i\alpha_{\nu_i,i}}\bigotimes_iG^{\nu_i}(G^iX),
\]
where $\Delta$ is the corresponding component of the exponential image of the additive diagonal. Write $\epsilon_X:B_1X\to X$ for projection onto degree one.

For $X=\bigoplus_{j=1}^sX_j$, label a source exponential summand by its multidegree $\lambda\in\mathbb N^s$. A target summand is indexed by multiplicities $\nu_\beta$ of summands of multidegree $\beta$ in $B_mX$; its total multidegree is $w=\sum_\beta\nu_\beta\beta$. The phrase \emph{multigrading preservation} below requires every component from $\lambda$ to this target to vanish unless $\lambda=w$, for every finite list of objects and every bound.

\begin{theorem}[Composition coherence]\label{prop:coherence}
Suppose the $D_{d,m}$ commute with projections in the bounds and satisfy the counit equations
\[
 \epsilon_{B_mX}D_{1,m}=\id,
 \qquad B_d(\epsilon_X)D_{d,1}=\id
\]
and the finite coassociativity equations
\[
 B_a(D_{b,c})D_{a,bc}=D_{a,b}^{B_cX}D_{ab,c}.
\]
Then bounded representatives $f:B_mX\to Y$ form a category, with representatives of higher bound identified by projection, identity $\epsilon_X$, and composition
\[
 g*f=g\,B_d(f)D_{d,m}\quad(g:B_dY\to Z).
\]
Conversely, if this formula is unital, associative and independent of the chosen degree bounds for all objects and representatives, the stated counit, coassociativity and projection equations hold. Compatibility with prescribed mixed-degree representations additionally requires multigrading preservation and, for $q$ objects $X_1,\ldots,X_q$ and $p\ge0$, the normalization
\[
 \operatorname{pr}_{(1,\ldots,1)}G^q((G^p\operatorname{pr}_j)_j)
 \alpha_{q,p}^{\oplus_jX_j}
 =\operatorname{pr}_{(p,\ldots,p)}.
\]
With these equations, the homogeneous and mixed representation maps are exactly the prescribed ones.
\end{theorem}

Here is the exactness statement. Let
\[
 \mathcal P_+(R)=\operatorname{Ch}^b_{\ge0}(\operatorname{Free}_{\mathrm{fg}}R)
\]
with degreewise split short exact sequences. Its disk $D_i$ has $R$ in homological degrees $i+1,i$ and identity differential. The map $u_i:D_i\to D_{i+1}$ is identity from the upper copy to the lower copy. An additive $R$-linear functor to an additive $R$-linear category $\mathcal B$ is called relatively left exact when it sends every distinguished short exact sequence to a kernel diagram.

\begin{theorem}[The all-ring essential image]\label{thm:walking}
For every commutative ring $R$, restriction $F\mapsto(F(D_i),F(u_i))$ is fully faithful on relatively left-exact functors. Its image consists exactly of the nonnegative cochain complexes $C$ in $\mathcal B$ for which, for every $L\in\mathcal P_+(R)$, the kernel of
\begin{align*}
 \delta_L(C):\bigoplus_i L_i\otimes_R C^i&\longrightarrow\bigoplus_iL_i\otimes_R C^{i+1},\\
 (x_i)_i&\longmapsto\bigl((1\otimes d^i)x_i-(\partial_{i+1}\otimes1)x_{i+1}\bigr)_i
\end{align*}
exists. All sums are finite. The inverse is $T_C(L)=\ker\delta_L(C)$. If $\mathcal B$ has all kernels, every such complex is in the image.
\end{theorem}
Full faithfulness and the case in which $\mathcal B$ has all kernels are already proved in Hua's exposition of the manuscript \cite[pp.~165--166]{Hua}. The criterion here specifies exactly which kernels suffice when arbitrary kernels are unavailable. This statement includes rings with zero divisors and the zero ring. It concerns ordinary kernels; requiring admissible kernels in a specified exact structure adds a separate condition. It does not require the nonnegative source category to be Frobenius.

Next, let $\mathcal T$ be the simplicial abelian groups whose normalized complex is bounded and finite free. For a simplicial set $X$ put
\[
 F_X(M)=\Hom_{s\mathsf{Ab}}(\Z[X],M).
\]
Evaluation defines natural linear maps $\rho_{p,M}:\Gamma^pF_X(M)\to F_X(\Gamma^pM)$.
\begin{theorem}[The full divided-power action]\label{thm:reconstruction}
On the image of $X\mapsto F_X$, the additive test functor with all its natural maps $\rho_{p,M}$ determines, and is determined by, the cosimplicial ring
\[
 A_X^n=\operatorname{Map}(X_n,\Z).
\]
Already $\rho_2$ and $\rho_0$ recover its tensor pairings and unit. This statement concerns the entire natural action on all test objects; operations on cohomology alone are not the asserted data.
\end{theorem}

The finiteness question has two different answers. Let $L$ be a
simplicial module over a commutative ring $R$ and let $N=N(L)$ be
its nonnegative normalized homological complex. A complex is
\emph{perfect} if it is quasi-isomorphic to a bounded complex of
finite projective modules, and \emph{pseudo-coherent} if it has a
finite-free resolution bounded below in homological degrees, possibly
unbounded above. The following criterion separates boundedness of
the given normalization from perfection of its derived class.
\begin{theorem}[Uniform rank bounds and derived fibres]\label{thm:rank-perfect}
\begin{enumerate}
\item Suppose every $L_n$ is finitely generated, and fix $D\ge0$.
Then $N_i=0$ for all $i>D$ if and only if, for every prime
$\mathfrak p$, the function
$n\mapsto\dim_{\kappa(\mathfrak p)}(L_n\otimes_R\kappa(\mathfrak p))$
is eventually a polynomial of degree at most $D$. The threshold may
depend on $\mathfrak p$, but $D$ is common. Maximal ideals suffice.
When the condition holds the exact formula, valid for every $n$, is
\[
 \dim_{\kappa(\mathfrak p)}(L_n\otimes_R\kappa(\mathfrak p))
 =\sum_{i=0}^D
 \dim_{\kappa(\mathfrak p)}(N_i\otimes_R\kappa(\mathfrak p))\binom ni.
\]
\item Suppose $R$ is noetherian and $N$ is pseudo-coherent.
Then $N$ is perfect if and only if there is a common $D\ge0$ with
\[
 H_i(N\otimes_R^{\mathbf L}\kappa(\mathfrak p))=0
 \quad\text{for every prime }\mathfrak p\text{ and every }i>D.
\]
Again maximal ideals suffice. This assertion uses derived fibres;
it does not assert a bound on the terms of the given $N$.
\end{enumerate}
\end{theorem}
The second assertion is the standard pseudo-coherent fibre criterion
\cite[Lemma~15.79.2, Tag~068V]{StacksFibres}. Section~\ref{sec:perfection}
gives the rank proof, the local mechanism of that criterion, and
counterexamples to replacing its derived fibres by ordinary ranks.

In cote~138, archive sheet~128 (PDF page~129), Grothendieck asks ``Unicit\'e d'un tel $\varphi$?'' for a clean holomorphic map with fixed embedded cells. The capstone about dessins answers the associated cellwise automorphism question. It uses the ordinary map obtained by suppressing the degree-two white vertices. A \emph{cellwise} automorphism fixes each black vertex, each resulting edge and each face setwise; it need not fix the white midpoint of an edge and need not commute with the Belyi map. The complete classification is the following.

\begin{theorem}\label{thm:dessins}
Let $X$ be a compact Riemann surface with a finite clean dessin. For holomorphic cellwise automorphisms of its ordinary map, the nontrivial groups are exactly these:
\begin{enumerate}
\item For the one-vertex one-loop sphere, in the model $K=\R\cup\{\infty\}$ with vertex $\infty$, the group is $z\mapsto az+b$ with $a>0$ and $b\in\R$.
\item For the sphere with vertices $0,\infty$ and $n\ge1$ equally spaced rays as edges, the group is $z\mapsto az$ with $a>0$.
\item For each $g\ge1$, the map with darts $\Z/(4g)$, vertex rotation $x\mapsto x+1$ and edge reversal $x\mapsto x+2g$ has group $C_2$.
\end{enumerate}
Every other connected finite clean map has trivial group. On a disconnected map the group is the product over components.
\end{theorem}

Mulase and Penkava classify ribbon graphs whose automorphism action on
edges has a nontrivial kernel \cite[Definition~1.9 and classification,
pp.~880--881]{MulasePenkava}. Requiring vertices and faces to be fixed
selects their one-vertex family with an even number of edges, which
is case~(3). That combinatorial classification is known. Here we
identify the full holomorphic group, whose kernel on darts need not
be trivial on the sphere, and include vertices of valence one and two.
Section~\ref{sec:dessins} gives this analytic argument and explicit
Belyi functions for every case.

The full cochain reconstruction has a precise homotopy consequence.
A space is of finite type here when its integral homology is finitely
generated in each degree; nilpotence means that its fundamental group
is nilpotent and acts nilpotently on all higher homotopy groups.

\begin{corollary}[Horel's and Mandell's recovery theorems applied to the action]
\label{cor:homotopy-recovery}
For connected nilpotent finite-type simplicial sets $X,Y$, the full
natural divided-power test action of Theorem~\ref{thm:reconstruction}
recovers their integral weak homotopy types and determines
\[
 \mathbb R\Map_{\mathsf{cBRing}}(A_Y,A_X)
 \simeq\mathbb R\Map_{\mathsf{Spaces}}(X,Y),
\]
where $A_X^n=\operatorname{Map}(X_n,\Z)$ has its unique pointwise
binomial operations, and $\mathsf{cBRing}$ is the category of
cosimplicial binomial rings with quasi-isomorphisms as weak
equivalences. Object detection also holds for a finite disjoint union
of such spaces. In the derived category of integral $E_\infty$
cochain algebras, object detection and faithfulness still hold, but
fullness fails: for $X=S^2,Y=S^3$, the cochain homotopy set is
$\widehat\Z$ whereas $[S^2,S^3]$ is a singleton.
\end{corollary}
This combines reconstruction with Horel's Theorem~5.4, p.~13, and
Mandell's Main Theorem and Theorems~0.1--0.2, pp.~213--216
\cite{Horel,Mandell}. Section~\ref{sec:homotopy-recovery} verifies
the binomial structure and the derived adjunction explicitly.

The dessin calculations lead to elementary arithmetic examples.
The classical sextic passport $[3,2,1;2,2,1,1;6]$ is Schneps's
three-member orbit \cite[Section~IV, Example~3, pp.~27--28]{Schneps}.
Section~\ref{sec:sextic} gives its complete normalized functions,
monodromy, fields, and collision parameters. The following theorem
specifies the related radical actions and the ambient groups in
which their images are open.

\begin{theorem}[Exact rank-one images and the scalar Kummer bound]
\label{thm:kummer-main}
Choose compatible roots $\theta_N=2^{1/N}>0$ and
$\zeta_N=e^{2\pi i/N}$. Write an element of the absolute Galois
group of $\Q$ as $(a,b)$ on them, so
$\sigma(\zeta_N)=\zeta_N^a$ and
$\sigma(\theta_N)=\zeta_N^b\theta_N$ at every $N$.
The image in $\widehat\Z\rtimes\widehat\Z^\times$, acting by
$x\mapsto ax+b$, is exactly
\[
 (-1)^b=\chi_2(a),\qquad
 \chi_2(a)=
 \begin{cases}1&a\equiv1,7\pmod8,\\-1&a\equiv3,5\pmod8.
 \end{cases}
\]
For the compatible roots $r_N=\zeta_{2N}\theta_N^{-1}$ of $-1/2$,
the translation is $c=(a-1)/2-b$, and its image is exactly
\[
 (-1)^c=\chi_{-2}(a),\qquad
 \chi_{-2}(a)=
 \begin{cases}1&a\equiv1,3\pmod8,\\-1&a\equiv5,7\pmod8.
 \end{cases}
\]
Each image has adelic index two and finite-level order
$N\varphi(N)/d_N$, where $d_N=2$ if $8\mid N$ and $d_N=1$
otherwise.

More generally, let $K$ be a number field and let
$\xi_1,\ldots,\xi_r\in K^\times$ be independent modulo roots of
unity. Choose a finite set $S$ of finite places containing every
nonzero valuation of the $\xi_j$. In
$U_S=\mathcal O_{K,S}^{\times}/\mu_K$, let $M$ be their span and
$M^{\mathrm{sat}}=(M\otimes\Q)\cap U_S$. If a positive integer $e$ annihilates
$M^{\mathrm{sat}}/M$, put $D=e|\mu_K|$.
The translation image over $K(\mu_\infty)$ contains
$D\widehat\Z^r$, and the image over $K$ in the scalar affine group
has index at most
\[
 [K\cap\Q(\mu_\infty):\Q]D^r.
\]
For $r\ge2$ that image has infinite index in the affine group with
full linear part $\operatorname{GL}_r(\widehat\Z)$.
\end{theorem}
The general openness is classical Kummer theory. Schinzel's abelian
radical obstruction, stated in \cite[Theorem~15, p.~7]{PeruccaSgobba},
and Perucca--Sgobba's uniform degree bounds
\cite[Theorems~1 and~13, pp.~1 and~6--8]{PeruccaSgobba}
are the relevant earlier results. Section~\ref{sec:kummer} proves
the explicit images and the displayed saturation bound, and derives
the diagonal action and indices for three Fermat fibres.


Classical divided powers already carry full composition relations.
For reduced operads over a field, Ikonicoff describes their algebras
by polynomial operations with an explicit substitution axiom
\cite[Definition~3.1.4 and Theorem~3.1.8, pp.~12--13]{Ikonicoff}.
His reduced hypothesis excludes nullary operations. The countermodels
here concern the constant and mixed-degree comparisons of the
manuscript's abstract tensor data.
Richter transports the simplicial tensor structure across Dold--Kan
and obtains divided-power chain algebras
\cite[Theorems~7.1 and~7.5, pp.~248--250]{Richter}.
Theorem~\ref{thm:reconstruction} instead specifies the natural action
on finite test objects that recovers the cosimplicial ring itself.

Recent formal work on polynomial laws by Chambert-Loir and de Frutos-Fern\'andez \cite[p.~75]{CLF} treats the classical universal constructions. The counterexamples above concern the extra abstract coherence asserted in the manuscript.
Fantechi and Ricolfi's account of the support morphism to a symmetric
product uses the classical polynomial-law and divided-power machinery
\cite[Appendix~B, pp.~27--30 of the author version]{FantechiRicolfi};
this geometric use also requires naturality under scalar extension.

The exact-category argument below uses B\"uhler's embedding and
cancellation theorems; its explicit matrix criterion applies over every
commutative ring. The counterexample when pullback closure is omitted
explains why that exact-category hypothesis is needed.

\paragraph{Why the arguments work.}
The homogeneous axioms see nested partitions but can miss constants interacting with two different degrees. A scalar twist preserves every homogeneous associativity equation while changing that interaction. The walking-complex theorem instead uses a canonical two-term copresentation by disks, so its only existence requirement is an explicit kernel. For cochains, divided powers recover cross terms, hence multiplication on evaluation objects; evaluation objects recover all cosimplicial degrees. In the dessin problem, an automorphism fixing one dart is forced to fix every dart by connectedness, leaving only one possible nonidentity permutation.

\section{Two counterexamples to the bare equivalence}
We give full models for both failures in Theorem~\ref{thm:bare}.

For the first take $\mathcal C=\mathsf{Vect}_{\F_2}^{\mathrm{fd}}$ and $G^nV=\Lambda^nV$. Exterior powers are exponential, with their signs disappearing in characteristic two. For positive $p,q$, define $\alpha_{p,q}$ by summing over unordered partitions of $pq$ input labels into $p$ blocks of size $q$, taking the exterior product within each block and then among the resulting blocks. This expression is multilinear and alternating: when two vectors coincide, terms in which their labels lie in one block vanish, and the other terms pair under transposition of those labels and cancel. For $q=1$ the repeated outer factors give zero directly. Thus it descends to $\Lambda^{pq}V$ and is natural.

Both sides of \eqref{eq:homogeneous} enumerate the same nested partitions into $p$ groups, each with $q$ blocks of $r$ labels. Every nested partition appears once, proving the equation. The one-index maps are identities. Put $\alpha_{0,q}=\id$ and put $\alpha_{p,0}:k\to\Lambda^pk$ equal to the canonical map, which is identity for $p=0,1$ and zero otherwise. Every zero-index equation is an identity or has both sides zero. Scalar degrees hold for the two scalars of $\F_2$. But $G^2(k)=0$, so $\alpha_{2,0}$ is not invertible.

For the second model use standard divided powers over $\F_3$ with their standard exponential comparisons. Set $t_2=-1$ and $t_n=1$ for $n\ne2$, including $0,1$. Replace the standard substitution maps by
\[
 \alpha'_{p,q}=c_{p,q}\alpha_{p,q},\qquad
 c_{p,q}=\frac{t_{pq}}{t_p t_q^p}.
\]
Cancellation gives $c_{p,qr}c_{q,r}^p=c_{pq,r}c_{p,q}$, so homogeneous associativity remains true. All one-index units and $\alpha'_{0,q}$ are unchanged, and all $\alpha'_{q,0}$ are invertible. Yet $\alpha'_{2,0}=-\alpha_{2,0}$.

Suppose these structured data came from polynomial composition with the fixed mixed and tensor-unit comparisons. Let $e:0\to k$ be the constant one and let $m:k\oplus k\rightsquigarrow k$ be multiplication, represented by the mixed component of $\Gamma^2(k\oplus k)$. Mixed-degree compatibility gives $m(e,\id)=\id$, hence $m(e,e)=e$. The homogeneous degree-two formula with the degree-zero pair instead gives
\[
 m(e,e)=m\Gamma^2((1,1))\alpha'_{2,0}=-e.
\]
This is impossible over $\F_3$. Rescaling homogeneous representatives cannot remove the discrepancy while preserving the specified exponential maps: a monoidal rescaling would require $t_{a+b}=t_at_b$, whereas $t_2\ne t_1^2$. This proves the theorem for the structured data in question.

\section{A finite coherence criterion for composition}\label{sec:coherence}
We prove the finite coherence criterion of Theorem~\ref{prop:coherence}, with the maps $B_m$, $D_{d,m}$ and $\epsilon$ defined in Section~1.

\begin{proof}
Projection compatibility makes composition independent of increasing a degree bound. The two counit equations give the left and right identities. For a third representative $h:B_aZ\to W$, expand both parenthesizations using the formula. Naturality of $D$ moves the linear map $f$ through the occurrence of $D$; the remaining equality is precisely the displayed coassociativity equation. This proves sufficiency.

For necessity choose the universal representatives given by identity maps on $B_mX$, then on $B_d(B_mX)$, and then on $B_aB_dB_mX$. Their identity and associativity equations read exactly as the equations for $D$; applying bound independence to these same universal representatives gives the projection equations. Increasing a bound is injective on representatives, since the projection onto the old bound has a section; equality after increasing bounds therefore gives these literal equations. Equality of maps can also be tested by all target maps, giving the Yoneda formulation. For a finite direct sum of source objects, expand both source and target into exponential summands. A target indexed by $(\nu_\beta)_\beta$ has multidegree $w_j=\sum_\beta\nu_\beta\beta_j$. Requiring the source component of multidegree $\lambda$ to vanish unless $\lambda=w$ is exactly multigrading preservation. The final equation says that the component representing a $q$-linear map in $q$ inputs of degree $p$ is its stipulated mixed component, with all other components zero. For inputs of unequal degrees, group the inputs of each degree $p$ in the exponential expansion and apply this normalization to that group. Coassociativity of the exponential diagonal combines the resulting groups, while multigrading eliminates the other components. Thus these equations are necessary and sufficient for the specified representation maps. Constants use $p=0$ and empty tensor products use $q=0$, so the missing mixed-unit case is included. All sums and tests are finite.
\end{proof}

This criterion records coherence rather than asserting that the bare homogeneous diagram implies it. Over a finite field the vanishing of off-diagonal multigraded components must be imposed explicitly; scalar eigenvalues alone do not separate all degrees.

\begin{counterexample}[Scalar degrees do not provide strict polynomial enrichment]\label{ex:enrichment}
Even with all unit maps invertible, the bare axioms do not force the maps
$f\mapsto G^p(f)$ on Hom modules to be homogeneous polynomial laws of
degree $p$, natural for composition.
\end{counterexample}
\begin{proof}
Work with finite-dimensional nonnegatively internally graded
$\F_2$-vector spaces and degree-zero maps. Let $T$ discard internal
degree zero and let $D^a$ be the ordinary divided power. Define the
external grading by
\[
 G^nX=\bigoplus_{a+3b=n}D^aX\otimes D^b(TX).
\]
The two gradings have different roles: internal degree is part of the
object, while $n$ labels the functor. Additivity of $T$ and the
exponential decomposition of divided powers give coherent exponential
comparisons for $G$. Moreover $G^0=\mathbf1$, $G^1=\id$ and
$G^n\mathbf1=\mathbf1$ for every $n$. The scalar-degree axiom holds
for both scalars in $\F_2$.

Let $i_n:D^n\to G^n$ and $\pi_n:G^n\to D^n$ be inclusion and
projection of the summand $b=0$. For $p,q\ge2$ define
\[
 \alpha^G_{p,q}
 =i_p^{G^qX}D^p(i_q^X)\alpha^D_{p,q}\pi_{pq}^X.
\]
Use identity maps for an index equal to one and the canonical unit
maps for a zero index. For three indices at least two, expand either
side of homogeneous associativity. Naturality moves each projection
past the adjacent divided-power map, and each $\pi_j i_j$ cancels.
Both sides become the standard threefold divided-power substitution,
preceded by $\pi_{pqr}$ and followed by the same three inclusions.
An index equal to one contributes an identity. If an index is zero,
the remaining objects are divided powers of the unit and all maps
are the identity of the unit. Thus these data satisfy every bare axiom.

Now place $V$ in internal degree one. The internal-degree-one summand
of $G^3V=D^3V\oplus V$ is $V$, and $G^3(f)$ restricts to $f$ on
that summand. A natural cubic enrichment would therefore give a
homogeneous cubic law $\ell$ on matrices agreeing with the identity
at $\F_2$-points. On one-dimensional spaces the only such law is
$r\mapsto r^3$: evaluation at $1$ fixes its coefficient. Coordinate
inclusions and projections force entrywise cubing on every matrix.
Postcomposition with $(x,y)\mapsto x+y$ then requires
$(x+y)^3=x^3+y^3$ in every $\F_2$-algebra. In $\F_2[s,t]$ their
difference is the nonzero polynomial $s^2t+st^2$. This contradiction
proves the claim. It also has a finite witness: in
$\F_4=\F_2[\omega]/(\omega^2+\omega+1)$, the two sides at
$(1,\omega)$ are respectively $1$ and $0$.
\end{proof}

\Needspace{7\baselineskip}
\subsection{The positive model on arbitrary modules}\label{sec:positive-laws}
The usual divided powers do satisfy the mixed identities, including for
modules that are neither flat nor finitely generated. We spell out this
positive model to distinguish it from the abstract counterexamples.

Roby's representing theorem gives
\[
 \operatorname{Pol}^d_k(X,Y)=\Hom_k(\Gamma^d_kX,Y)
\]
for arbitrary modules over a commutative ring. His scalar-extension and
exponential theorems give, respectively,
\[
 R\otimes_k\Gamma_k^dX\simeq\Gamma_R^d(R\otimes_kX),\qquad
 \Gamma^d(X\oplus Y)\simeq
 \bigoplus_{a+b=d}\Gamma^aX\otimes_k\Gamma^bY.
\]
These are \cite[Theorem~IV.1, pp.~266--267; Theorems~III.3--III.4,
pp.~262--263]{Roby}. In particular, invariant tensors must not replace
the universal divided powers on an arbitrary nonflat module.

\begin{proposition}\label{prop:ordinary-polynomial-category}
The modules over $k$, with morphisms
$\bigoplus_{i\ge0}\Hom_k(\Gamma^iX,Y)$, form the category of
polynomial laws with a globally finite degree bound. Composition is
composition after every scalar extension. If
$f=\sum_{i=0}^m f_i$ and $g_q:\Gamma^qY\to Z$, its degree-$n$ part is
\begin{equation}\label{eq:ordinary-mixed}
 \sum_{\substack{\sum_i\nu_i=q\\\sum_i i\nu_i=n}}
 g_q\mu_\nu
 \left(\bigotimes_i\Gamma^{\nu_i}(f_i)\right)
 \left(\bigotimes_i\alpha_{\nu_i,i}\right)
 \Delta_{(i\nu_i)_i}.
\end{equation}
Here $\alpha_{a,b}:\Gamma^{ab}X\to\Gamma^a\Gamma^bX$
represents $x\mapsto\gamma_a(\gamma_b(x))$, $\Delta$ is the
specified component of the divided-power diagonal, and $\mu_\nu$
is multiplication into $\Gamma^qY$. The homogeneous degree-one subcategory is
$\mathsf{Mod}_k$, and finite products are direct sums.
\end{proposition}
\begin{proof}
For a homogeneous pair, naturality of the universal law gives
$g_q\Gamma^q(f_p)\alpha_{q,p}$. For a sum of homogeneous laws,
apply the divided-power identity
\[
 \gamma_q(y_0+\cdots+y_m)=
 \sum_{\nu_0+\cdots+\nu_m=q}
 \gamma_{\nu_0}(y_0)\cdots\gamma_{\nu_m}(y_m)
\]
to $y_i=f_i(\gamma_i(x))$. Its term indexed by $\nu$ has input
degree $\sum i\nu_i$. Representing that term gives exactly
\eqref{eq:ordinary-mixed}. The sum is finite; the constant term
causes no exception, since $\alpha_{q,0}(1)=\gamma_q(1)$.
Composition is associative and unital as an operation on natural
families of functions. The representing bijection transfers these
identities to the displayed formula. A pair of bounded laws has a
common finite bound and a unique law to the direct sum. Degree-one
laws are precisely linear maps, again by representation. The
exponential decomposition represents a law of multidegree
$(d_1,\ldots,d_s)$ on $X_1\oplus\cdots\oplus X_s$ by
$\bigotimes_j\Gamma^{d_j}X_j$. Thus the representing tensor
products are the usual tensor products of modules, with their
ordinary associativity, permutation and unit maps. Uniqueness of
all degree decompositions follows by substituting $tx$ over $R[t]$
and comparing powers of the indeterminate $t$.
\end{proof}

\begin{proposition}\label{prop:gamma-hom-law}
For arbitrary $k$-modules $X,Y$ and $p\ge0$, the operation
\[
 f\longmapsto\Gamma^p(f):
 \Hom_k(X,Y)\rightsquigarrow
 \Hom_k(\Gamma^pX,\Gamma^pY)
\]
is a homogeneous polynomial law of degree $p$. It preserves identities
and composition as identities of laws and is natural in both modules.
No base-change isomorphism for $\Hom$ is assumed.
\end{proposition}
\begin{proof}
For finitely many maps $f_1,\ldots,f_s$, expand
$\Gamma^p(\sum_iT_if_i)$ over $k[T_1,\ldots,T_s]$.
Its coefficient of $T^\nu$, for $\sum\nu_i=p$, is
\[
 c_\nu=\mu_\nu\left(\bigotimes_i\Gamma^{\nu_i}(f_i)\right)
          \Delta_\nu
 \in\Hom_k(\Gamma^pX,\Gamma^pY).
\]
For $\sum_i r_i\otimes f_i\in R\otimes_k\Hom_k(X,Y)$,
define the value to be $\sum_\nu r^\nu\otimes c_\nu$.
We must check that it depends only on this tensor, even when the
canonical map from the tensor product of a Hom module to a Hom module
after base change is not an isomorphism.

Adjoining a zero term, splitting either a map or a coefficient as
a sum, and balancing a $k$-scalar between its coefficient and the map generate all changes
of tensor presentation. The divided-power sum and scalar identities
show that the indicated coefficient operators respect these changes.
For the splitting of a coefficient on a repeated map, the needed
identity is $\mu_{a,b}\Delta_{a,b}=\binom{a+b}{a}\id$ on
$\Gamma^{a+b}$; it follows by expanding $\gamma_{a+b}((s+t)x)$
over $k[s,t]$ and comparing coefficients.
One can check each identity with independent indeterminates: a finite
polynomial of $k$-linear operators is zero after extension to a free
polynomial algebra only when each coefficient operator is zero.
Indeed, evaluate on any input and compare its coefficients in the
free module of monomials. Thus the coefficient equalities hold in
the ordinary Hom module itself and can be specialized to every $R$.
This proves independence of presentation and naturality in $R$.
Every monomial has degree $p$, so the law is homogeneous.
The ordinary functorial identities for $\Gamma^p$ over the free
polynomial algebra give identity and composition compatibility by
the same coefficient argument. For $p=0$ the resulting law is the
constant identity of $k$.
\end{proof}

\begin{proposition}\label{prop:three-polynomial-categories}
Over a nonzero commutative ring $k$, the category of bounded polynomial
laws, the fixed-degree category with morphisms
$\Gamma^d\Hom_k(V,W)$ on finite projective modules, and its category
of linear representations must be distinguished. The latter
representations are exactly functors whose action on linear Hom
modules is a degree-$d$ polynomial law compatible with composition
and identities. The bounded polynomial-law category is equivalent
to neither of the other two categories.
\end{proposition}
\begin{proof}
For $k=\F_2$, $V=k$, $W=k^2$ and $d=2$, the module
$\Hom_k(\Gamma^2V,W)$ has dimension two, whereas
$\Gamma^2\Hom_k(V,W)$ has dimension three. Thus even the proposed
Hom modules differ. Homogeneous polynomial laws of degree $d>1$
compose in degree $d^2$, and have no degree-one identities within
that single piece.

There is also a categorical obstruction valid for every nonzero
$k$. The bounded-law category has terminal object zero and no
initial object: every object has at least the two distinct constant
laws $0,1$ to $k$. For $d>0$, the fixed-degree category has a zero
object. For $d=0$, each of its Hom sets is $k$, so it has no
terminal object. Its representation category is abelian and has a
zero object for every $d$. These properties are invariant under
equivalence, proving the non-equivalence assertions.

Finally Roby's representation applied to the module
$\Hom_k(V,W)$ identifies a degree-$d$ law with a linear map
\[
 \Gamma^d\Hom_k(V,W)\longrightarrow\Hom_k(FV,FW).
\]
The composition and identity equations represent exactly the
functor equations in the fixed-degree category. The same argument
identifies the natural transformations. This is the standard strict
polynomial-functor category of \cite[Section~2, pp.~3--5]{Krause}.
\end{proof}


\subsection{Exactly what scalar weights can distinguish}
The scalar-degree caveat has a precise arithmetic form. For an abelian
 group $M$ write $M[d]=\{x:dx=0\}$.
\begin{proposition}[Scalar-weight ambiguity]\label{prop:scalar-weights}
Let $A=\bigoplus_{m\ge0}A_m$, with the scalar $\lambda\in\Z$ acting
as $\lambda^m$ on $A_m$, and use $0^0=1$. Put
\[
 d_{n,m}=\gcd_{\lambda\in\Z}|\lambda^n-\lambda^m|\quad(n\ne m).
\]
The simultaneous weight-$n$ eigenspace is exactly
\[
 A_n\oplus\bigoplus_{m\ne n}A_m[d_{n,m}].
\]
If one index is zero then $d_{n,m}=1$. For positive distinct indices
put $s=\min(n,m)$ and $h=|n-m|$. For an odd prime $\ell$,
\[
 v_\ell(d_{n,m})=
 \begin{cases}
 0,&\ell-1\nmid h,\\
 \min\{s,1+v_\ell(h)\},&\ell-1\mid h,
 \end{cases}
\]
and
\[
 v_2(d_{n,m})=
 \begin{cases}
 \min\{s,1\},&h\text{ odd},\\
 \min\{s,2+v_2(h)\},&h\text{ even}.
 \end{cases}
\]
Over a commutative ring $k$, replace $(d_{n,m})$ by the ideal
$I_{n,m}(k)=(\lambda^n-\lambda^m:\lambda\in k)$. For several
scalar variables the corresponding ideal is
$\sum_j I_{n_j,m_j}(k)$. Distinct positive weights over $\F_q$
coincide on every scalar exactly when $q-1$ divides their difference;
over an infinite field every pair of distinct weights is separated.
Naturality after all commutative scalar extensions separates distinct
weights over every ring, with no torsion hypothesis.
\end{proposition}
\begin{proof}
Write $x=\sum_m x_m$. The weight equation holds exactly when
$(\lambda^m-\lambda^n)x_m=0$ for every $m,\lambda$. The ideal of
these integers is generated by their gcd, giving the eigenspace
formula by B\'ezout. If one exponent is zero, evaluation at zero gives
one. Otherwise the difference, up to sign, is
$\lambda^s(\lambda^h-1)$.

For a prime $\ell$, scalars divisible by $\ell$ have minimum valuation
$s$, attained at $\lambda=\ell$. For odd $\ell$, the units modulo
$\ell$ form a cyclic group of order $\ell-1$. If this order does
not divide $h$, a unit gives valuation zero. If it does, write
$h=(\ell-1)r$. Every unit satisfies
$v_\ell(\lambda^h-1)\ge1+v_\ell(r)$: expand a power of
$1+\ell a$ successively, first by an exponent prime to $\ell$ and
then by powers of $\ell$. Equality is attained. Indeed, lift a
primitive nonzero residue to an integer $u$ with
$v_\ell(u^{\ell-1}-1)=1$. Such a lift exists because replacing $u$
by $u+t\ell$ changes the next coefficient by the nonzero residue
$(\ell-1)u^{\ell-2}t$. The same binomial expansion gives valuation
$1+v_\ell(r)=1+v_\ell(h)$ for this lift. Taking the smaller of
the unit and nonunit minima proves the odd-prime formula.

For an odd integer $u$ and odd $h$, $u^h-1$ has valuation
$v_2(u-1)$, whose minimum is one, attained at $u=3$. If
$h=2^r h'$ with $h'$ odd and $r\ge1$, successive factorization of
$u^{2^r}-1$ gives
\[
 v_2(u^h-1)=v_2(u-1)+v_2(u+1)+r-1.
\]
The first two valuations sum to at least three, with equality at
$u=3$, so the unit minimum is $2+v_2(h)$. Comparison with $s$
proves the remaining formula.

The argument over $k$ is the same componentwise ideal-annihilator
calculation. For multidegrees, setting all but one scalar equal to
one gives each coordinate ideal; a telescoping expansion of the
difference of two products gives the reverse inclusion. The finite
field assertion follows from the exponent of $\F_q^\times$, and
the infinite field assertion from the fact that a nonzero polynomial
has only finitely many roots. Finally take the scalar $t$ in $k[t]$.
For $n\ne m$, the two different coefficients of
$(t^n-t^m)x$ force $x=0$ in $A\otimes_k k[t]$. This works even
when $A$ has torsion or $k$ has zero divisors.
\end{proof}
The formula describes torsion ambiguity, not a demand to invert
primes: a torsion-free abelian group already has no off-weight
vectors. It also explains why scalar identities at the ground-ring
points cannot substitute for polynomial-law naturality.

\section{The walking complex over an arbitrary ring}\label{sec:walking}
We prove Theorem~\ref{thm:walking} by a canonical copresentation. For a finite free module $V$, $V\otimes_R B$ means a finite copower in $\mathcal B$; a change of basis acts by its matrix, so the construction is independent of the basis.

For $b\in\mathcal B$, set $C_b^i=\Hom_{\mathcal B}(b,C^i)$ and let $L^\vee$ be the cochain complex dual to $L$. The kernel formula gives naturally
\[
 \Hom_{\mathcal B}(b,T_C(L))
 =\Hom_{\operatorname{Ch}(R)}(L^\vee,C_b).
\]
The right side denotes ordinary cochain maps, not derived maps. It proves functoriality. Dualizing a degreewise split exact sequence reverses it, and Hom into $C_b$ is left exact. Testing on every $b$ therefore proves relative left exactness of $T_C$. A map from $D_i^\vee$ is determined by its degree-$i$ component, giving $T_C(D_i)=C^i$ and $T_C(u_i)=d^i$.

Conversely define
\[
 I(L)=\bigoplus_{i\ge0}D_i\otimes_RL_i.
\]
In degree $n$ this has $L_n\oplus L_{n-1}$, with differential $(x,y)\mapsto(y,0)$. The map $j_L(x)=(x,\partial x)$ is a degreewise split monomorphism. Its quotient is $Q(L)_n=L_{n-1}$ with differential $-\partial$, and its quotient map is $q_L(x,y)=y-\partial x$. Both
\[
 0\to L\xrightarrow{j_L}I(L)\xrightarrow{q_L}Q(L)\to0
\]
and the analogous sequence for $Q(L)$ are distinguished. Relative left exactness makes $F(L)$ the kernel of
\[
 F(I(L))\xrightarrow{F(j_{Q(L)}q_L)}F(I(Q(L))).
\]
Additivity identifies this with $\delta_L(C)$ for $C=\rho F$: next-disk entries are $d^i$ and same-disk entries are $-\partial_{i+1}$. Thus $F\simeq T_{\rho F}$. A cochain map induces maps on these kernels; conversely every natural transformation is determined by its values on the disk copresentations. This proves full faithfulness and the essential-image statement.

\begin{proposition}[Enough projectives and injectives, with a boundary]
\label{prop:nonnegative}
The category $\mathcal P_+(R)$ has enough projectives and injectives.
For $R\ne0$ it is not Frobenius: the stalk $S_0$ in degree zero is
projective and is not injective.
\end{proposition}
\begin{proof}
A disk is projective because maps from $D_i\otimes V$ are determined
by maps from $V$ to degree $i+1$; it is injective because maps into
it are determined by maps from degree $i$ to $V$. Here $V$ is
finite free and all sequences split degreewise. The copresentation
$I(L)$ above therefore supplies enough injectives. For projectives,
use
\[
 J(L)=(S_0\otimes L_0)\oplus
       \bigoplus_{i\ge1}(D_{i-1}\otimes L_i).
\]
In degree $n$ the map to $L$ is
$(x,y)\in L_n\oplus L_{n+1}\mapsto x+\partial y$.
It is a degreewise split chain epimorphism, and its kernel has finite
free terms isomorphic to $L_{n+1}$. The stalk $S_0$ is projective
since $\Hom(S_0,L)=L_0$, so $J(L)$ is projective. Finally the
inclusion $S_0\to D_0$ has no chain retraction: a map $D_0\to S_0$
must vanish in degree zero by its chain-map equation. For $R\ne0$
this proves the failure of injectivity and of the Frobenius property.
\end{proof}

\subsection{What changes for a principal ideal domain}
The arbitrary-ring criterion simplifies over a principal ideal domain,
but it does not reduce to the kernels of the individual differentials.
\begin{proposition}\label{prop:pid}
Let $R$ be a principal ideal domain. In Theorem~\ref{thm:walking} it
suffices, and is necessary, that every $\ker d^i$ and every kernel
\[
 K_{i,a}=\ker\left[C^i\oplus C^{i+1}\longrightarrow
 C^{i+1}\oplus C^{i+2},\quad
 (x,y)\longmapsto(d^ix-ay,d^{i+1}y)\right]
\]
exist, for all nonzero nonunits $a\in R$. Over a field the second
family is empty. For every nonfield principal ideal domain the
first family alone is insufficient in general.
\end{proposition}
\begin{proof}
For a bounded finite-free complex $L$, its cycles $Z_i$ and boundaries
$B_i$ are finite free. Split $L_i\to B_{i-1}$, and put each inclusion
$B_i\subset Z_i$ into Smith normal form. The resulting bases express
$L$ as a finite direct sum of stalks and two-term complexes
$R\xrightarrow{a}R$. A unit $a$ gives a disk. The matrix kernel of
a stalk is $\ker d^i$, that of a disk is $C^i$, and that of the
displayed two-term complex is $K_{i,a}$. Kernels of finite direct
sums are their direct sums, proving the equivalence.

For insufficiency choose a prime element $\pi$ of $R$ and let
$\mathcal B$ be the category of finite free $S=R/(\pi^2)$-modules.
Take $C^1=S$ and all other terms zero. All differential kernels
exist in $\mathcal B$, but $K_{0,\pi}$ would represent the kernel
of multiplication by $\pi$ on $S$. Evaluation of Hom at $S$ would
identify its underlying $S$-module with $\pi S\simeq R/(\pi)$.
That nonzero module is killed by $\pi$ and is not free over $S$.
Thus the required kernel does not exist in $\mathcal B$.
\end{proof}

\subsection{Exact categories and admissible kernels}
For a Quillen exact target the exact-functor statement has a useful
form in terms of its admissible short exact sequences, or
\emph{conflations}. The first nonzero arrow of a conflation is an
\emph{inflation}, and the second is a \emph{deflation}.
\begin{theorem}\label{thm:exact}
Let $\mathcal E$ be an $R$-linear exact category. Restriction to the
disks is fully faithful on exact functors from $\mathcal P_+(R)$.
Its essential image consists of complexes admitting objects $Z_i$
and conflations
\[
 0\longrightarrow Z_i\xrightarrow{s_i}C^i
 \xrightarrow{p_i}Z_{i+1}\longrightarrow0,
 \qquad d^i=s_{i+1}p_i \quad(i\ge0).
\]
In an abelian target this means $H^i(C)=0$ for $i>0$; $H^0(C)$
is unrestricted. For relatively left-exact functors additionally
required to preserve inflations, the exact criterion is existence
of the kernels of Theorem~\ref{thm:walking} and admissibility of
every induced map $T_C(A)\to T_C(B)$ for a source inflation
$A\to B$. If $\mathcal E$ is weakly idempotent complete, it is
enough to require that every canonical kernel inclusion
$T_C(L)\to\bigoplus_i L_i\otimes_R C^i$ be an inflation.
\end{theorem}
\begin{proof}
The source conflation $0\to S_i\to D_i\to S_{i+1}\to0$
shows necessity, with $Z_i=F(S_i)$. Full faithfulness follows from
Theorem~\ref{thm:walking}.

For sufficiency first work in an abelian target and form the internal
Hom complex $\mathcal H_L=\operatorname{Hom}^{\bullet}(L^\vee,C)$,
so that $\mathcal H_L^q=\bigoplus_i L_i\otimes_R C^{i+q}$ and
$Z^0\mathcal H_L=T_C(L)$. Filtering $L^\vee$ by its finitely many
degrees gives successive internal-Hom quotients that are finite
sums of $C[i]$ with $i\ge0$. Their cohomology in degree $q>0$
is zero. The long exact sequence of each filtration step gives
$H^q(\mathcal H_L)=0$ for $q>0$. A source conflation gives a
degreewise split short exact sequence of these Hom complexes.
The obstruction to lifting a zero-cycle lies in
$H^1(\mathcal H_A)=0$, so the sequence of zero-cycles is short
exact. This proves exactness of $T_C$ in the abelian case.

Use the exact fully faithful extension-closed abelian embedding of
an exact category \cite[Theorem~A.1, p.~36]{Buehler}. Work in a universe
containing the objects in question; the embedding preserves their
central $R$-actions. It remains to show that each abelian kernel
$T_C(L)$ lies in $\mathcal E$. Let $p$ be the least nonzero degree
of $L$, put $V=L_p$, and let $L'$ be its quotient by the stalk
$S_p\otimes_R V$. In complexes, $L$ is the pullback of
$D_p\otimes_RV\to S_{p+1}\otimes_RV$ along the map from $L'$
induced by $\partial_{p+1}$. This follows degree by degree: the
only new coordinate is $V$ in degree $p$, and the compatibility
equation in degree $p+1$ is exactly that differential. Applying
the abelian construction identifies $T_C(L)$ with the pullback of
the deflation $C^p\otimes_RV\to Z_{p+1}\otimes_RV$ along
$T_C(L')\to Z_{p+1}\otimes_RV$. The Quillen pullback axiom and
induction on the support of $L$ put this object in $\mathcal E$.
Reflection of exactness by the embedding completes the proof.

For the last assertion let $k_L$ denote the canonical kernel
inclusion. Necessity follows by applying the functor to $j_L$.
Conversely, for an inflation $a:A\to B$ in the source, $I(a)$
is split since its degreewise maps split. The equation
$k_BT_C(a)=T_C(I(a))k_A$ says that this composite is an inflation.
Cancellation of inflations in a weakly idempotent complete exact
category \cite[dual of Proposition~7.5, p.~17]{Buehler} makes $T_C(a)$
an inflation. This proves the shorter test.
\end{proof}

The two elementary closure axioms on sequences do not suffice in
place of Quillen exactness. Take $R=\Z$. In finite abelian groups let $\Sigma$
consist of finite direct sums of split sequences and copies of
\[
 e:\quad0\longrightarrow\Z/4\xrightarrow{4}\Z/16
 \longrightarrow\Z/4\longrightarrow0.
\]
This is a class of kernel--cokernel pairs closed under isomorphism
and finite direct sum and containing every split sequence. The
complex $C^0=\Z/16\to C^1=\Z/4$ given by reduction, with all
other terms zero, is a resolution for this class. Any exact extension
of its disk values would, by Theorem~\ref{thm:walking}, have to be
$T_C$. Applied to $0\to S_0\to[R\xrightarrow{2}R]\to S_1\to0$,
it has middle group
\[
 G=\{(x,y)\in\Z/16\oplus\Z/4:x\bmod4=2y\bmod4\}.
\]
This group has order sixteen and exponent eight: its first
coordinates are even and $(2,1)$ has order eight. Its kernel and
quotient under projection to $y$ both have order four and are
cyclic. A sequence in $\Sigma$ with those two ends is either
split, with middle exponent four, or is a copy of $e$, with middle
exponent sixteen. Neither is possible. This explicitly isolates
the missing pullback closure in the weaker axioms.

\subsection{The two-sided Frobenius replacement}
The obstruction at degree zero disappears when complexes may extend
in both directions.
\begin{proposition}\label{prop:frobenius}
For every commutative ring $R$, give
$\operatorname{Ch}^b(\operatorname{Free}_{\mathrm{fg}}R)$ in all
integer degrees its degreewise split exact structure. It is a
Frobenius category: its projective and injective objects coincide
and are exactly the contractible complexes. Its quotient by maps
factoring through these objects is
$K^b(\operatorname{Free}_{\mathrm{fg}}R)$, the category of bounded
finite-free complexes and chain-homotopy classes of maps.
\end{proposition}
\begin{proof}
A disk in degrees $i+1,i$ with coefficients in a finite free module
$V$ is projective and injective for this exact structure. Maps out
of it are determined by a map from $V$ to degree $i+1$; maps into
it are determined by a map from degree $i$ to $V$. Degreewise
splittings lift or extend these maps.

For a complex $C$, its cone on the identity is
\[
 P(C)_n=C_{n-1}\oplus C_n,\qquad
 d(a,b)=(-d_Ca,a+d_Cb).
\]
The map $H(a,b)=(b,0)$ contracts it. There is a degreewise split
sequence $0\to C\to P(C)\to C[1]\to0$; shifting it also gives
a deflation from a contractible complex to $C$. Every contractible
complex is projective and injective. To see this without assuming
its cycles are free, use the disk sum
$I(C)_n=C_n\oplus C_{n-1}$, with differential $(x,y)\mapsto(y,0)$.
The usual inclusion $c\mapsto(c,d_Cc)$ has, for a contracting
homotopy $h$ of $C$, the chain retraction
\[
 (x,y)\longmapsto d_Chx+hy.
\]
Thus $C$ is a retract of a finite sum of disks, and inherits their
relative projectivity and injectivity. The cone sequences give enough
of both. Conversely, an injective splits its inclusion in $P(C)$,
and a projective splits the shifted cone deflation. Either is a
retract of a contractible complex and hence contractible. This proves
the Frobenius assertion, including rings with zero divisors.

A map factoring through a contractible complex is null-homotopic,
by composing its contracting homotopy with the two maps. Conversely,
if $f:C\to D$ has $f=d_Dh+hd_C$, it factors through the cone above:
\[
 C\longrightarrow P(C)\longrightarrow D,\qquad
 c\longmapsto(0,c),\quad (a,b)\longmapsto h(a)+f(b).
\]
The chain-map equation for the second arrow is exactly the displayed
null-homotopy equation. Thus the ideal being quotiented out is
precisely the chain-null-homotopic maps, proving the stable-quotient
statement.
\end{proof}

\section{Reconstruction from the natural divided-power action}\label{sec:reconstruction}
We construct both directions in Theorem~\ref{thm:reconstruction}. The category $\mathcal T$ is closed under degreewise tensor products and divided powers: Dold--Kan writes the rank in simplicial degree $m$ as a polynomial combination of $\binom mi$, and tensor products multiply these polynomials while $\Gamma^p$ replaces rank $r$ by $\binom{r+p-1}p$. Binomial inversion makes the normalized ranks zero in sufficiently high degree. Normalized groups are direct summands of finite free abelian groups and hence free.

The exponential decomposition
\[
 \Gamma^2(P\oplus Q)=\Gamma^2P\oplus(P\otimes Q)\oplus\Gamma^2Q
\]
recovers a tensor pairing from $\rho_2$: include the middle summand of $\Gamma^2(F_X(M)\oplus F_X(N))$, apply $\rho_{2,M\oplus N}$, and project to $F_X(M\otimes N)$. Evaluation at a simplex gives exactly $f(x)\otimes g(x)$. The map $\rho_0:\Z\to F_X(\Z)$ gives the unit.

Define $R_n$ in simplicial degree $m$ by
\[
 (R_n)_m=\Z^{\Hom_\Delta([n],[m])}.
\]
Its simplicial maps restrict functions along composition. Dold--Kan and Vandermonde's identity give
\[
 \operatorname{rank}(NR_n)_i=\binom{n+1}i\ (0\le i\le n+1),
 \qquad (NR_n)_i=0\ (i>n+1).
\]
Thus $R_n\in\mathcal T$. Its normalization is in fact split
acyclic. On the unnormalized chain complex define
$h_m:(R_n)_m\to(R_n)_{m+1}$ by extending a function to be zero on
maps $[n]\to[m+1]$ whose image contains $0$, and otherwise shifting
all their values down by one. Then $d_0h=\id$ and
$d_jh=hd_{j-1}$ for $j>0$, with $d_1h_0=0$ at the bottom.
The alternating differential therefore satisfies $\partial h+h\partial=\id$.
The normalized and degenerate complexes split naturally, and the
degenerate complex is acyclic, so the normalization is acyclic as
well. An acyclic bounded complex of finite free abelian groups splits,
since its cycle groups are free and each sequence onto the next
cycle group splits. This proves the split-acyclic assertion without
using it to erase the nontrivial evaluation functor.

The adjunction with degree-$n$ evaluation is explicit: a function $f:X_n\to\Z$ sends $x\in X_m$ to the function $\theta\mapsto f(X(\theta)x)$, and the inverse evaluates at $\id_{[n]}$. Hence $F_X(R_n)=\Z^{X_n}$. The maps $R_n\to R_{n'}$ recover every cosimplicial operator. Pointwise multiplication $R_n\otimes R_n\to R_n$, followed after the recovered tensor pairing, gives pointwise multiplication on $\Z^{X_n}$. This recovers the entire cosimplicial ring.

Conversely, since $M_m$ is finite free,
\[
 A_X^m\otimes M_m\simeq\operatorname{Map}(X_m,M_m)
\]
without any assumption on the cardinality of $X_m$. The functor $F_X(M)$ is the equalizer of
\[
 \prod_m A_X^m\otimes M_m
 \rightrightarrows\prod_{\theta:[m]\to[n]}A_X^n\otimes M_m,
\]
where the two arrows apply $A_X(\theta)$ and $M(\theta)$ respectively. Multiplication of the $A_X$ coefficients recovers tensor pairings. In a basis $e_1,\ldots,e_r$ of $M_m$, divided powers are recovered by
\[
 \gamma_p\left(\sum_jf_je_j\right)
 =\sum_{\alpha_1+\cdots+\alpha_r=p}
   \left(\prod_jf_j^{\alpha_j}\right)
   \gamma_{\alpha_1}(e_1)\cdots\gamma_{\alpha_r}(e_r).
\]
The divided-power identities give basis independence and compatibility with simplicial operators; universal linearization gives $\rho_p$. All identities can be checked by evaluation at each simplex, so the two reconstructions are inverse. Normalizing the underlying cosimplicial abelian group gives the integral cochain complex. Cosimplicial Dold--Kan recovers that underlying group from the complex; its multiplication is retained by the full action above. This proves the stated reconstruction.

\subsection{What the reconstructed ring says about homotopy}\label{sec:homotopy-recovery}
The reconstruction theorem identifies the algebraic input to existing
integral homotopy-recovery theorems. Here we state exactly the consequences
and distinguish the two categories in which derived maps are taken.

\begin{proof}[Proof of Corollary~\ref{cor:homotopy-recovery}]
For each $n$, the function ring $A_X^n=\operatorname{Map}(X_n,\Z)$
is torsion-free and has pointwise binomial operations
$f\mapsto\binom f r$. These operations are uniquely determined by
\[
 r!\binom f r=f(f-1)\cdots(f-r+1),
\]
because multiplication by $r!$ is injective. Every ring map between
such rings preserves them. Consequently the cosimplicial ring of
Theorem~\ref{thm:reconstruction} canonically determines its
cosimplicial binomial-ring structure.

Horel's model category of cosimplicial binomial rings has weak
equivalences the quasi-isomorphisms of normalized cochain complexes
\cite[Theorem~3.7, p.~7]{Horel}. His adjunction sends a simplicial
set $X$ to $\Z^X$ and has derived right adjoint
\[
 \langle A\rangle=\mathbb R\Map_{\mathsf{cBRing}}(A,\Z)
 \qquad\text{\cite[Section~4, p.~9]{Horel}}.
\]
For connected nilpotent finite-type $Y$, his Theorem~5.4, p.~13,
says that the unit $Y\to\langle A_Y\rangle$ is a weak equivalence.
The derived adjunction therefore gives
\[
 \mathbb R\Map_{\mathsf{cBRing}}(A_Y,A_X)
 \simeq\mathbb R\Map_{\mathsf{Spaces}}(X,\langle A_Y\rangle)
 \simeq\mathbb R\Map_{\mathsf{Spaces}}(X,Y).
\]
These are applications of Horel's theorem, rather than a new
homotopy-recovery theorem.

For the integral $E_\infty$ cochain category, Mandell's Main Theorem
\cite[p.~213]{Mandell} says that the cochain algebra detects weak
homotopy types of finite-type nilpotent spaces. His Theorem~0.1,
p.~214, gives faithfulness. His Theorem~0.2, p.~215, describes the
extra maps: the right adjoint applied to $C^*(Y;\Z)$ is the homotopy
pullback of the free loop space of the finite completion $Y^\wedge$
along $Y\to Y^\wedge$. Thus for simply connected $Y$ its $n$th
homotopy group is
\[
 \pi_nY\oplus\widehat{\pi_{n+1}Y}.
\]
In particular, for $X=S^2$ and $Y=S^3$, the homotopy set $[X,Y]$
is a singleton, whereas the set of derived $E_\infty$ cochain maps
is $\widehat\Z$ and is uncountable; this is Mandell's explicit
example on p.~216. The two mapping-space assertions refer to
different derived categories.

Finally let $X$ have finitely many components, each nilpotent and
of finite type. The primitive idempotents of
$H^0(A_X)=\Z^{\pi_0X}$ recover the component set. The normalized
complex is nonnegative, so a degree-zero cohomology class has a
unique cocycle representative. For each primitive idempotent this
representative is the characteristic function of one component on
vertices. Its cocycle condition makes it constant on the endpoints
of every edge. Consequently all its cosimplicial images in degree
$n$ agree on the vertices of each simplex and give the characteristic
idempotent of that component of $X_n$. These idempotents split
$A_X$ canonically, up to permutation, into the function rings of its
components. Apply the connected object-detection theorem to each
factor. The zero function ring distinguishes the empty space.
\end{proof}


To\"en's integral affinization of a simply connected finite-type
space gives another comparison. Its homotopy sheaves are the homotopy
groups tensored with the Hilbert additive group scheme, and its
integral points have a natural retraction onto the original space
\cite[Theorem~4.1, p.~11; Corollaries~5.4--5.5, pp.~16--17]{Toen}.
The unit to those points need not be an equivalence: Corollary~5.1,
p.~14, gives additional completed homotopy groups, including
$\widehat\Z$ in degree two for $S^3$. The binomial category in
Corollary~\ref{cor:homotopy-recovery} has the stated unit equivalence.


\subsection{What weaker cochain data fail to recover}
\begin{counterexample}\label{ex:lens}
The connected nilpotent finite CW spaces $L(5,1)$ and $L(5,2)$ have
quasi-isomorphic integral cochain complexes and isomorphic graded
cohomology rings, including every degree-compatible divided power on
the positive even cohomology ideal. They are not homotopy equivalent.
\end{counterexample}
\begin{proof}
Use the convention
\[
 L(5,q)=S^3/\bigl((z,w)\sim(\zeta_5z,\zeta_5^qw)\bigr),
 \qquad q=1,2.
\]
The action is free, so the fundamental group is $C_5$ and the
universal cover is $S^3$. Every deck transformation is connected to
the identity through coordinate rotations of $S^3$; it acts trivially
on the higher homotopy groups. Thus these spaces are nilpotent.
Their integral cohomology is
\[
 H^0=H^3=\Z,\qquad H^2=\Z/5,\qquad H^1=0,
\]
with no other groups. One obtains this from $H_1=C_5$, orientation,
Poincar\'e duality and the universal coefficient theorem. All products
of positive-degree classes vanish for degree reasons. On the positive
even ideal $H^2$, every divided power of order at least two also
vanishes for degree reasons. These are the same operations on both
spaces; this assertion does not put the unit in a divided-power ideal.

A finite cellular cochain complex over $\Z$ splits, by Smith normal
form, into stalks, two-term multiplication blocks and contractible
disks. The groups just listed therefore make either complex
quasi-isomorphic to
\[
 \Z\xrightarrow{0}\Z\xrightarrow{5}\Z\xrightarrow{0}\Z
 \quad\text{in cohomological degrees }0,1,2,3.
\]
This is a statement about complexes, without the full natural
multiplicative action of Theorem~\ref{thm:reconstruction}.

For the specified quotient orientation, the linking form is
$\lambda_q(a,b)=-qab/5$ in $\Q/\Z$
\cite[Proposition~1.2, p.~2]{CFH}. A homotopy equivalence carries
this form to itself up to the degree $\pm1$: the defining Bockstein,
cup product and evaluation on the fundamental class are natural.
An isomorphism on $H_1=\Z/5$ is multiplication by a unit $u$.
Consequently an equivalence between these two spaces would require
$2\equiv\pm u^2\pmod5$. The nonzero squares modulo five are
$1,4$, and changing the sign gives the same set. This is impossible.
\end{proof}

\begin{counterexample}[Binomial operations are a different structure]\label{ex:binomial}
The operation $a\mapsto\binom a2$ on $\Z$ is not the point map
of a polynomial law from the rank-one integral module to itself.
Moreover, no divided-power ideal of $\Z$ contains $1$.
\end{counterexample}
\begin{proof}
A law would be natural for the reduction $\Z\to\F_2$. The integers
$0$ and $2$ have the same image there, whereas their proposed outputs
are $0$ and $1$. Naturality is contradicted. Equivalently, its
universal polynomial would be $t(t-1)/2$, which is not in $\Z[t]$.
If a divided-power ideal contained $1$, the identity
$1^2=2\gamma_2(1)$ would make $1$ divisible by $2$ in $\Z$.
\end{proof}
Thus binomial operations in homotopy-recovery theorems, an internal
divided-power ideal, and the natural action on the entire test
functor are distinct structures. The reconstruction theorem uses the
last of these.

\section{Simplicial ranks and derived perfection}\label{sec:perfection}
Let $L$ be a simplicial module over a commutative ring $R$, with
normalized homological complex $N=N(L)$ in degrees $i\ge0$.
We call this particular normalized complex \emph{strictly perfect}
when its terms are finite projective and zero in sufficiently high
degree. We call it \emph{perfect} when it is quasi-isomorphic to a
bounded complex of finite projective modules. The latter condition
allows a different representative and ignores contractible summands.

\subsection{The rank formula and its converse}
The Dold--Kan decomposition has the concrete form
\begin{equation}\label{eq:dold-kan-sum}
 L_n\simeq\bigoplus_{i=0}^n N_i^{\oplus\binom ni}.
\end{equation}
The summands are indexed by increasing surjections $[n]\twoheadrightarrow[i]$;
there are $\binom ni$ such maps, since their $i$ jumps can occur at
any $i$ of the $n$ gaps. This is the denormalization formula in
Dold--Puppe \cite[equation~(3.3), p.~219; Theorem~3.6 and
Corollary~3.8, p.~220]{DoldPuppe}. In particular, $N_i$ is a natural
direct summand of $L_i$. It follows in both directions that all
$L_n$ are finitely generated, respectively finite projective, exactly
when all $N_i$ have that property. The decomposition is compatible
with arbitrary scalar extension: its splitting maps are made from
faces and degeneracies. This also proves that normalization commutes
with scalar extension, despite its possible description by kernels.

For a prime $\mathfrak p$ write
\[
 a_i(\mathfrak p)=\dim_{\kappa(\mathfrak p)}
       (N_i\otimes_R\kappa(\mathfrak p)),\qquad
 r_{\mathfrak p}(n)=\dim_{\kappa(\mathfrak p)}
       (L_n\otimes_R\kappa(\mathfrak p)).
\]
When the terms are finite, \eqref{eq:dold-kan-sum} gives the exact,
nonnegative binomial expansion
\begin{equation}\label{eq:fibre-rank}
 r_{\mathfrak p}(n)=\sum_{i=0}^n a_i(\mathfrak p)\binom ni.
\end{equation}
\begin{proof}[Proof of the rank assertion in Theorem~\ref{thm:rank-perfect}]
If $N_i=0$ for $i>D$, the displayed formula is a polynomial in $n$
of degree at most $D$ for every prime, valid from $n=0$. Conversely,
suppose the fibre ranks are eventually polynomials of degree at most
$D$, with the threshold and coefficients allowed to depend on the
prime. If $a_j(\mathfrak p)>0$ for some $j>D$, nonnegativity gives
$r_{\mathfrak p}(n)\ge\binom nj$ for every $n\ge j$, contradicting
that polynomial bound as $n$ tends to infinity. Thus all fibres of
$N_j$ vanish for $j>D$. Each $N_j$ is finitely generated, so
Nakayama's lemma at the maximal ideals implies $N_j=0$. Maximal
ideals therefore suffice for the test. For a nonzero fibre the degree
of its polynomial is the largest $i$ with $a_i(\mathfrak p)>0$,
since its leading coefficient is positive. The zero polynomial has
all coefficients zero. The statement also covers the zero ring.
\end{proof}
The common degree bound is essential; the starting index of eventual
polynomial behavior need not be common. If the simplicial terms are
finite projective, this theorem characterizes strict perfection.
Without projectivity, it characterizes bounded normalization only.

\subsection{Perfect complexes over regular rings}
\begin{proposition}\label{prop:regular-perfect}
Let $R$ be a regular noetherian ring, without a bound on its Krull
dimension. Then $N(L)$ is perfect if and only if every $H_i(N(L))$
is finitely generated and these groups vanish in sufficiently high
degree. In particular, finite simplicial terms and the common degree
bound in Theorem~\ref{thm:rank-perfect} imply perfection.
\end{proposition}
\begin{proof}
This is the regular-ring perfectness criterion
\cite[Lemma~15.76.14, Tag~066Z]{StacksRegular}. Its finite-global-
bound point can be seen directly. For a finitely generated module
$M$, choose a resolution with finite free terms; noetherianity keeps
its syzygies finite. The loci where successive syzygies are locally
free are increasing open sets. Each prime belongs to one of them,
since its regular local ring has finite global dimension. Quasi-
compactness of $\Spec R$ gives a single syzygy that is locally free
everywhere, and hence finite projective. Thus this particular $M$
has a finite projective resolution, even when there is no uniform
bound valid for all finite modules.

A complex with only finitely many nonzero finite homology modules
is built in the derived category from their shifts by its finite
truncation triangles. Each of these modules has just been shown
perfect; taking a finite number of mapping cones preserves perfectness.
Conversely, a bounded complex of finite projectives has bounded
finitely generated homology over a noetherian ring. Finally, finite
simplicial terms and a common rank degree bound make $N$ bounded
with finite terms by \eqref{eq:dold-kan-sum} and the rank assertion;
its homology has the required finiteness.
\end{proof}

\begin{counterexample}[Two different failures of a rank test]\label{ex:rank-tests}
Over $\Z$, fibrewise eventual polynomial ranks with no common degree
bound do not imply perfection. Perfection does not imply eventual
polynomial ranks of a particular termwise finite simplicial model.
\end{counterexample}
\begin{proof}
For the first assertion choose distinct primes $p_i$, $i\ge1$, and
take the zero-differential complex $N_i=\Z/p_i$, $N_0=0$.
Denormalize it. Every simplicial degree is a finite abelian group.
At the prime $p_j$ its rank function is $\binom nj$, and at the
generic point it is zero. Thus every fibre is eventually polynomial,
but the degrees are unbounded. Its homology is nonzero in infinitely
many degrees, so it cannot be perfect.

For the second take the locally finite direct sum of the disks with
$\Z$ in degrees $i+1,i$, for all $i\ge0$. Then $N_0$ has rank one
and every $N_i$ for $i>0$ has rank two. The disk contractions sum
degreewise, giving a contracting homotopy. Hence $N$ represents zero
and is perfect. Its simplicial rank in degree $n$ is nevertheless
\[
 1+2\sum_{i=1}^n\binom ni=2^{n+1}-1,
\]
which is not eventually polynomial. Each simplicial term is still
finite free. The same example works over every nonzero ring.
\end{proof}

\subsection{The exact test uses derived fibres}
Recall that a complex is \emph{pseudo-coherent} if it has a resolution
by finite free modules bounded below in homological degrees, though
possibly unbounded above. For a nonnegative complex over a noetherian
ring, finite terms, or more generally finite homology in every degree,
suffice: choose finitely many cycle representatives in the bottom
degree, take a finite free module on them, and continue through the
successive cones. At each step noetherianity makes the needed kernels
and homology modules finite. This constructs the finite free resolution.

\begin{proof}[Proof of the derived-fibre assertion in Theorem~\ref{thm:rank-perfect}]
The exact external criterion is
\cite[Lemma~15.79.2, Tag~068V]{StacksFibres}: a pseudo-coherent
complex has finite projective amplitude in a fixed interval if and
only if its derived residue-field fibres vanish outside that interval;
maximal ideals suffice. Here $N$ is nonnegative, so its derived
fibres have no negative homological degrees. The asserted common
upper bound therefore gives the interval $[0,D]$.

For completeness, the local mechanism behind the criterion in this
case is elementary. Choose a finite-free resolution $P$ in nonnegative
homological degrees, not necessarily bounded above. At a local ring,
a unit entry in a differential splits off a contractible rank-one
disk by row and column operations. Starting in the bottom degree,
repeat this operation until all remaining differential entries are
in the maximal ideal. Only finitely many operations affect each
degree. The remaining minimal complex has zero differential after
tensoring with the residue field. Vanishing of the derived fibre
above $D$ therefore makes all its free terms above $D$ zero.
Consequently $P$ is locally homotopy equivalent to a complex supported
in $[0,D]$, and
$\operatorname{coker}(P_{D+1}\to P_D)$ is locally finite free.
This cokernel is finitely presented, so it is finite projective
when the condition holds at all maximal ideals. Also $H_i(P)=0$
for $i>D$, as can be tested by localization. Replacing degree $D$
by that cokernel and removing all higher terms gives a bounded
projective complex quasi-isomorphic to $P$. This proves sufficiency.
Necessity follows by tensoring any bounded projective representative.
The use of a common degree interval and of derived tensor products is
essential.
\end{proof}

Even bounded normalization with constant ordinary fibre rank can fail.
Let $R=\F_2[\epsilon]/(\epsilon^2)$, let $M=R/(\epsilon)$, and let
$L$ be the constant simplicial module $M$. Its normalization is $M$
in degree zero, and its ordinary residue-field rank is one in every
simplicial degree. But
\[
 \cdots\xrightarrow{\epsilon}R\xrightarrow{\epsilon}R
 \xrightarrow{\epsilon}R\longrightarrow M\longrightarrow0
\]
is a free resolution: the kernel and image of multiplication by
$\epsilon$ are both $(\epsilon)$. Tensoring with $M$ makes every
differential zero, so
$\operatorname{Tor}^R_i(M,M)=M$ for every $i\ge0$. The derived
fibres have no common upper bound, and $M$ is not perfect.

\subsection{Constructible coefficients and the PID block formula}
\begin{proposition}\label{prop:constructible-pid}
If $R$ is noetherian and each $L_n$ is finite, every function
$\mathfrak p\mapsto a_i(\mathfrak p)$ in
\eqref{eq:fibre-rank} is constructible. If $R$ is a principal ideal
domain and every normalized term is finite free, write $\mu_i$ for
the number of unit differential blocks in degrees $i,i-1$, write
$\beta_{i,a}$ for the number of blocks $R\xrightarrow{a}R$ there,
with $a$ a nonzero nonunit up to associates, and write $\nu_i$ for
the number of stalks in degree $i$. Put
$\beta_i=\sum_a\beta_{i,a}$ and $\mu_0=\beta_0=0$. Then
\begin{equation}\label{eq:pid-ranks}
 \operatorname{rank}_R N_i
 =\mu_i+\mu_{i+1}+\beta_i+\beta_{i+1}+\nu_i.
\end{equation}
The actual normalized complex is strictly perfect exactly when the
three families $\mu,\beta,\nu$ have finite support. It is perfect
exactly when $\beta,\nu$ have finite support; no finite-support
condition on the contractible family $\mu$ is needed.
\end{proposition}
\begin{proof}
The module $N_i$ is finite by \eqref{eq:dold-kan-sum}, hence finitely
presented over the noetherian ring. In a presentation
$R^s\xrightarrow{A}R^r\to N_i\to0$, its fibre dimension is
$r-\operatorname{rank}A(\mathfrak p)$. The minors of $A$ describe
each rank stratum as a locally closed set. There are finitely many
possible ranks, proving constructibility.

Over a PID, the cycles and boundaries of a degreewise finite-free
complex are free. Split each surjection $N_i\to B_{i-1}$ and
apply Smith normal form to each inclusion $B_i\subset Z_i$.
As in Proposition~\ref{prop:pid}, this decomposes the complex into
the stated two-term blocks and stalks. It works in unbounded
nonnegative degrees as a locally finite sum: every term contains
only finitely many blocks. A block with upper degree $i$ contributes
to both degrees $i,i-1$, giving \eqref{eq:pid-ranks}; in particular
the term $\mu_i$ cannot be omitted.

Boundedness of the actual terms is exactly finite support of all
three families. Unit blocks are contractible. A nonunit block with
upper degree $i$ contributes $R/(a)$ to homology in degree $i-1$,
and a stalk contributes $R$ in its own degree. There is no
cancellation between these direct summands. Thus bounded homology is
exactly finite support of $\beta,\nu$, which is equivalent to
perfection over the regular noetherian PID by
Proposition~\ref{prop:regular-perfect}. This proves both criteria.
The block assertion requires free terms; it is not a classification
of arbitrary complexes with torsion terms over a PID.
\end{proof}

\section{Holomorphic groups preserving the cells of a clean dessin}\label{sec:dessins}
We now prove Theorem~\ref{thm:dessins}. Let $\beta:X\to\mathbf P^1$ be a nonconstant clean Belyi map on a connected compact Riemann surface, and suppress its white degree-two points as vertices. Let $G_{\rm cell}$ preserve each black vertex, edge and face, with the conventions of Section~1.


\begin{proof}
The combinatorial reduction is the vertex- and face-fixed
specialization of Mulase--Penkava's exceptional-graph classification
\cite[pp.~880--881]{MulasePenkava}. We give its dart form to include
the low-valence maps in the statement. Let $\sigma$ and $\alpha$ be
the vertex and edge permutations of the darts. Connectedness makes $\langle\sigma,\alpha\rangle$ transitive. An automorphism commuting with both and fixing one dart fixes every dart. If a nonidentity such permutation preserves each edge, it must therefore send every dart to its partner, so it equals $\alpha$. Preservation of each vertex and commutation with $\sigma$ force $\sigma$ to be one cycle. Write its length as $2m$; then $\alpha=\sigma^m$. The face permutation has $\gcd(2m,m-1)=\gcd(2,m-1)$ cycles. Its cycles are individually preserved by translation by $m$ exactly when $m$ is even. Thus $m=2g$, giving one vertex, $2g$ edges and one face, and genus $g$. Conversely these permutations give the stated cellwise involution.

Every combinatorial automorphism of a dessin is realized by a unique Belyi automorphism: over the complement of the three branch values, lift it by the monodromy covering; local inverse branches make the lift holomorphic, and the local models $z\mapsto z^e$ extend it over branch points. Let $H$ be the subgroup of $G_{\rm cell}$ fixing every dart. At a black vertex its derivative is positive real, because it fixes each outgoing analytic ray. In positive genus the automorphism group fixing a point is finite. A finite-order holomorphic germ with positive derivative has derivative one and is the identity: averaging a local coordinate over its cyclic group produces an invariant coordinate with nonzero derivative. The identity theorem gives $H=1$. The combinatorial calculation therefore proves the positive-genus assertion.

On the sphere every induced permutation of the darts is trivial. Three fixed vertices force a M\"obius transformation to be the identity. With exactly two vertices at $0,\infty$ an element is $z\mapsto az$ with $a>0$. If $a\ne1$, its iterates force every edge to join those vertices; invariance then forces every edge to be a ray: the iterates $a^{-k}z$ approach zero with constant argument, which must be the unique incident tangent direction (replace the element by its inverse if necessary). The local Belyi model makes the rays equally spaced. Conversely positive dilations preserve all rays and sectors.

With one vertex at infinity an element is affine, $az+b$, with $a>0$. If $a\ne1$, translating its second fixed point to zero shows that the closure of every invariant edge contains zero. As zero is not a vertex, disjoint edge interiors force there to be only one edge. If $a=1,b\ne0$, all invariant edges approach infinity in the same two tangent directions; distinct darts have distinct tangents, again forcing one edge. Its clean model is $\beta(z)=1/(1+z^2)$ with the real-circle graph. Preserving its two faces forces exactly the stated real affine group. This proves the list. Cellwise preservation forbids permutations of connected components, giving the product assertion.
\end{proof}

The second family has explicit clean Belyi function
\[
 \beta_n(z)=-\frac{4z^n}{(z^n-1)^2}.
\]
Its graph is the union of rays $z^n\in\R_{\le0}$, and $\beta_n-1=-(z^n+1)^2/(z^n-1)^2$ verifies cleanliness. For the third family use
\[
 X_g:\ y^2=x^{2g+1}-x,\qquad\beta_g=x^{2g}.
\]
The polynomial is square-free. The double cover has $2g+2$ branch points and genus $g$; the Belyi map has degree $4g$, one black vertex, one face, and $2g$ white points of index two. The transformation $(x,y)\mapsto(\zeta_{4g}^2x,\zeta_{4g}y)$ generates a cyclic deck group; its $2g$th power is $(x,y)\mapsto(x,-y)$, the cellwise involution.

\subsection{Counting maps with labelled cells}\label{sec:pinned-count}
An action on cells has both an image and a kernel, and the two play
different roles in a counting problem. Let an oriented map have finite
dart set $E$, vertex rotation $\rho$, and edge involution $\alpha$;
we use $\phi=\rho^{-1}\alpha$ for its face permutation. Fixed points
of $\alpha$ are allowed and describe half-edges. Connectedness means
transitivity of $\langle\rho,\alpha\rangle$.

The classical covering correspondence identifies transitive permutation
data with Belyi pairs and cycle lengths with ramification indices
\cite[Lemma~I.1, p.~2, and Theorem~I.5, p.~8]{Schneps}. A permutation
commuting with both generators gives a unique covering transformation
by path lifting. The local power models extend it over the branch
points. Thus the Belyi automorphism group is their simultaneous
centralizer. For degree $d$ the genus is determined by
\[
 \#\operatorname{cyc}(\rho)+\#\operatorname{cyc}(\alpha)
 +\#\operatorname{cyc}(\phi)=d+2-2g.
\]
In the clean case this becomes the usual vertices-minus-edges-plus-faces
formula, since $d$ is twice the number of ordinary edges.

Fix a type $T$ consisting of graded sets of vertices, edges and faces,
where the grade is the degree. Write
$\gamma_T=|\operatorname{Aut}(T)|$; if the counts in degree $j$ are
$a_j,b_j,c_j$, then $\gamma_T=\prod_j a_j!b_j!c_j!$.
A pinning labels the graded cells of a map by $T$. For an unlabelled
map $C$, let $A_C$ be its full combinatorial group, and let $K_C,I_C$
be the kernel and image of its action on those graded cells.

\begin{proposition}\label{prop:pinning-count}
If $\mathcal M(T)$ is the finite set of unpinned isomorphism classes
of type $T$, with any specified connectedness or genus restrictions,
then the number of pinned isomorphism classes is
\[
 N_T=\sum_{[C]\in\mathcal M(T)}\frac{\gamma_T}{|I_C|}
 =\gamma_T\sum_{[C]\in\mathcal M(T)}\frac{|K_C|}{|A_C|}.
\]
Each pinning of $C$ has automorphism group $K_C$. Equality with
$|\mathcal M(T)|\gamma_T$ holds exactly when every $I_C$ is trivial;
pinned rigidity instead means every $K_C$ is trivial.
\end{proposition}
\begin{proof}
For fixed $C$, its pinning set is a torsor under $\operatorname{Aut}(T)$.
The action of $A_C$ factors through $I_C$, whose action on this torsor
is free. Hence every orbit has size $|I_C|$, while every stabilizer
in $A_C$ is $K_C$. Sum the orbit counts over distinct unpinned maps.
Each summand is at most $\gamma_T$, with equality precisely when
$|I_C|=1$, proving also the asserted criterion. The argument includes
disconnected maps and the empty class set. If the empty map is allowed,
it has one pinning of empty type; a nonempty clean map has even dart
number.
\end{proof}

For the smallest counterexample to identifying these conditions, take
$E=\Z/4$, $\rho(x)=x+1$ and $\alpha(x)=x+2$.
Here $\phi=\rho$, so there is one vertex, two edges and one face,
on a torus. The centralizer is the translation group $C_4$; its even
translations fix the two edges and its odd translations interchange
them. Thus $K_C=C_2$, $I_C=C_2$ and $\gamma_T=2$.
There is only one unpinned map of this type: of the three pairings
of four cyclically ordered darts, only the opposite pairing gives
one face. The formula therefore gives one pinned class. The kernel
classification in the proof of Theorem~\ref{thm:dessins} identifies
this as the genus-one member of the complete one-vertex one-face
family, and shows that any half-edge forces $K_C=1$.

\subsection{Functions on fixed embedded cells and noncompact maps}
The group in Theorem~\ref{thm:dessins} preserves the reduced map,
whose white points have been suppressed. It differs from the group
commuting with its Belyi function. This distinction also matters for
noncompact dessins: Cañas, Hidalgo, Turiel and Viruel realize every
countable group as a dessin automorphism group
\cite[Theorems~1--2, p.~3; Section~2.3, p.~4]{Canas}.
Their automorphisms may permute cells. The group considered here
fixes every cell separately.

\begin{proposition}\label{prop:fixed-cell-torsor}
Fix a connected compact Riemann surface $X$ with a clean dessin's embedded
reduced cells, and let $H$ be the subgroup of its holomorphic
cellwise group fixing every dart. The clean Belyi functions having
exactly these cells form a free transitive $H$-set under precomposition.
Thus nonuniqueness occurs exactly for the one-loop sphere and the
oranges of every order. Marking every face centre removes these
continuous groups. Marking every white point removes them for
oranges but leaves positive dilations for the one-loop sphere.
\end{proposition}
\begin{proof}
Two such functions give identical oriented dart data. The covering
isomorphism specified by the identity of these data extends
holomorphically across each branch point, by the local power-map
argument above. It fixes each dart and intertwines the functions,
so belongs to $H$. Conversely precomposition by $H$ preserves the
embedded cells. A stabilizer would be a Belyi automorphism fixing a
dart, and connectedness then makes it the identity. This proves the
torsor assertion. The compact classification identifies $H$ with
the positive affine group or positive dilations in the two spherical
families, and with the trivial group in positive genus. The finite
positive-genus involution already preserves the Belyi function.

For $1/(1+z^2)$ the white point is $0$ and the two face centres are
$i,-i$. An affine map $az+b$, $a>0,b\in\R$, fixes $0$ precisely when
$b=0$, and fixes $i$ only when also $a=1$. In an orange every white
point and face centre is finite and nonzero, so fixing any one of them
forces a positive dilation to have scale one.
\end{proof}
The torus involution reverses the two loops of its one-vertex
one-face structure and acts as $-1$ on $H_1(-;\Z)$; it is therefore
not isotopic to the identity. Conversely a M\"obius transformation is
isotopic to the identity on the sphere but can move its dessin.
Cellwise preservation and isotopy to the identity are different conditions.

\begin{proposition}\label{prop:noncompact-cell-rigidity}
Let a connected noncompact Riemann surface carry a locally finite
clean holomorphic cellular map, with compact edge closures and
ordinary disk faces attached along their finite edge cycles.
Every cellwise holomorphic automorphism is the identity.
\end{proposition}
\begin{proof}
The transitive dart action has a centralizer acting freely on darts.
If a nonidentity permutation in it fixes each edge, it must reverse
every edge, hence equal $\alpha$. Fixing every vertex then forces
one vertex rotation cycle. Local finiteness makes that cycle finite,
so there are finitely many darts, edges and face cycles. Their
ordinary closed-cell attachments make the surface compact, a
contradiction. Thus any cellwise automorphism fixes all darts.

Its derivative at a black vertex is positive real because it
preserves each outgoing analytic ray. Uniformization
\cite[Theorem~4.2.1 and the quotient classification, pp.~152--153]{Wendl}
reduces to a disk or plane universal cover. In the disk case lift
the map to fix a chosen lift of the vertex. After moving it to zero,
Schwarz's lemma makes the lift a rotation; its positive derivative
forces the identity. For a plane cover with a nontrivial deck group,
noncompactness leaves a rank-one translation lattice. The lift fixing
a chosen vertex is affine, and its positive multiplier preserves
this lattice, so is one. Finally, on $\C$ a nonidentity positive
dilation about the fixed vertex has a two-sided orbit unbounded in
one direction. An interior point of an invariant compact edge
closure cannot have such an orbit. The map is therefore the identity
in every case.
\end{proof}
The ordinary-cell hypothesis excludes asymptotic tracts; an arbitrary
holomorphic map to the sphere need not have the required cellular structure.

\section{Schneps's sextic example and its normalization}\label{sec:sextic}
This section makes the known three-member sextic orbit explicit in the
normalization suggested by the manuscript. Schneps's cubic in
\cite[Section~IV, Example~3, pp.~27--28 of the author version]{Schneps}
is $25t^3-6t^2-6t-2$; setting $c=2t$ gives the cubic below.
The cubic and its three-member orbit are credited to that example.
The factorization here also separates the collision parameters from
the required passport.

Set
\begin{align*}
 h_c(z)&=z^3(z-1)^2(z-c),&
 G_c(z)&=z^2-\frac{5c+4}{6}z+\frac c2,\\
 \delta(c)&=25c^2-32c+16,&p(c)&=25c^3-12c^2-24c-16,\\
 P(c)&=-\frac{\delta(c)p(c)}{7776},&
 Q(c)&=\frac{c(5c-8)(25c^3-6c^2+8)}{2592}.
\end{align*}
Differentiation gives $h_c'=6z^2(z-1)G_c$. Division by $G_c$ gives
remainder $P(c)z+Q(c)$. To verify it directly, put
$\alpha=(5c+4)/6$, $\beta=-c/2$, so $z^2=\alpha z+\beta$
in the quotient. Reduce the six powers in $h_c$ using this relation;
the resulting coefficients are
\begin{align*}
 7776P&=-625c^5+1100c^4-184c^3-176c^2-128c+256,\\
 2592Q&=125c^5-230c^4+48c^3+40c^2-64c.
\end{align*}
Factoring them gives the displayed formulas. If $G_c$ has distinct
roots $u,v$, their critical values agree precisely when
$P(c)(u-v)=0$, hence, after excluding $\delta=0$, when $p(c)=0$.
The raw remainder coefficient has degree five; the effective
noncollision equation has degree three.

\begin{proposition}\label{prop:sextic-complete}
The three roots of $p$ give exactly the Belyi pairs with ordered
passport $[3,2,1;2,2,1,1;6]$ through
\[
 \beta_c(z)=\frac{1500z^3(z-1)^2(z-c)}{32-52c-c^2}.
\]
They have trivial Belyi automorphism groups. Writing
$r_j=\sqrt[3]{2}\,\zeta_3^j$ with positive real $\sqrt[3]{2}$,
the parameters are
$c_j=(4+18r_j+6r_j^2)/25$. Their fields of moduli and definition are
$\Q(c_j)=\Q(r_j)$; their splitting field is
$\Q(\sqrt[3]{2},\zeta_3)$, with Galois group $S_3$ and quadratic
subfield $\Q(\sqrt{-3})$.
\end{proposition}
\begin{proof}
The values of $p$ modulo $7$ at $0,\ldots,6$ are
$5,1,4,3,1,1,6$. Thus $p$ is irreducible over $\Q$.
Its discriminant is $-2^8 3^9$, so it has one real and two nonreal
roots, all distinct. Modulo $p$,
\[
 Q=\frac{32-52c-c^2}{1500},\qquad
 h_c-Q=G_c^2H_c,
 \quad H_c=z^2+\frac{2c-2}{3}z+\frac{5c^2-4c-4}{12}.
\]
For an explicit identity before reduction, the difference
$h_c-Q-G_c^2H_c$ is
\[
 \frac{p(c)}{2592}
 (-5c^2-4c-30cz^2+36cz+24z^3-24z^2).
\]
None of $c,1-c,Q,\delta$ vanishes at a root of the irreducible cubic.
The discriminant of $H_c$ is $(16+4c-11c^2)/9$, also nonzero.
The resultant of $G_c,H_c$, reduced modulo $p$, is
$(c-1)(5c^2+2c+2)/6$, so these factors have no common root.
Their zeros give exactly two double and two simple white vertices.
The derivative has no other finite critical points, and infinity is
a pole of order six. This proves the passport.

Conversely that passport has total ramification $3+2+5=10$;
Riemann--Hurwitz makes its source a sphere. Put its unique pole at
infinity, its triple black point at zero and its double black point
at one. The normalized function is uniquely $\lambda h_c$.
Its residual critical points are distinct and have value one, forcing
$p(c)=0$ and $\lambda=Q(c)^{-1}$. Every pair is therefore on the list.
An isomorphism fixes the three distinguished points and hence is the
identity; it also forces equality of the parameters.

Substitute $c=(4+18r+6r^2)/25$ and reduce by $r^3-2$ to obtain
$p(c)=0$. Both $\Q(c)$ and $\Q(r)$ have degree three and the former
is contained in the latter, so they coincide. The three substitutions
give all conjugates. The irreducible cubic's nonsquare discriminant
makes the Galois group $S_3$; its square root generates
$\Q(\sqrt{-3})$. Finally, Galois conjugation sends $\beta_c$ to
$\beta_{\sigma(c)}$. The unique normalization shows that these pairs
are isomorphic exactly when $\sigma(c)=c$. This identifies the field
of moduli directly, and the formula gives a model over it.
\end{proof}

The corresponding monodromy representatives, composing right to left,
are
\[
\begin{array}{c|c|c}
 &\alpha&\phi=\rho^{-1}\alpha\\\hline
 T_1&(14)(36)&(154362)\\
 T_0&(14)(56)&(156432)\\
 T_2&(14)(26)&(154326)
\end{array}
\qquad\rho=(123)(45)(6).
\]
Each face permutation is a six-cycle, so each map is connected and
has genus zero. The centralizer of $\rho$ is $C_3\times C_2$.
In $T_0$ the singleton dart $6$ is paired with a dart of the two-cycle
of $\rho$, whereas in $T_1,T_2$ it is paired with a dart of the
three-cycle. The first distinction is invariant. To carry $T_1$ to
$T_2$, a centralizing permutation would have to carry $3$ to $2$;
its unique rotation on $(123)$ then carries the other paired dart
$1$ to $3$, not to $1$, so it cannot carry $(14)$ to $(14)$.
Thus the three classes are distinct, and the geometric completeness
proof shows they exhaust the passport. Conjugation by $(23)$ reverses
$\rho$ and exchanges $T_1,T_2$, while fixing $T_0$.
The real parameter therefore gives $T_0$, and the nonreal parameters
give the mirror pair. A separate exact enumeration finds all eighteen
one-face involutions in three centralizer orbits of six.

For completeness, the roots of $\delta$ are
$16/25\pm12i/25$. The repeated residual critical point
$u=(5c+4)/12$ is distinct from $0,1,c$, and
$Q\equiv-2(139c-82)/9375\pmod\delta$ is nonzero.
The normalized polynomial then has white passport $[3,1,1,1]$,
because its derivative has order two at $u$ and no other residual
critical point. These parameters must be excluded from the preceding
proposition. The ancillary coefficient identities are
\[
 (\alpha-2)(\alpha-c)+\beta+1
 =-\frac{5c^2-10c-4}{36},\qquad
 \alpha(\alpha^2+2\beta)
 =\frac{(5c+4)(25c^2+4c+16)}{216}.
\]
The two double white points are the roots of $G_c$; their multiplicity-two factor in $f-1$ is $G_c^2$.
More generally, logarithmic differentiation of coprime factorizations
gives
\[
 \frac{f'}f=\sum_i\frac{p_i}{z-a_i}-\sum_j\frac{q_j}{z-b_j}
 \quad\text{for }f=\lambda
 \frac{\prod_i(z-a_i)^{p_i}}{\prod_j(z-b_j)^{q_j}}.
\]


\section{Exact Kummer images and diagonal Fermat fibres}\label{sec:kummer}
The Galois actions in the manuscript can be made explicit by choosing
compatible roots and keeping the cyclotomic multiplier in the affine
group. The choice of origin changes the translation coordinate, so
we specify it throughout.

Write $\widehat\Z=\varprojlim_N\Z/N$ and put
\[
 \operatorname{Aff}_1(\widehat\Z)
 =\widehat\Z\rtimes\widehat\Z^\times,
 \qquad (a,b)(a',b')=(aa',b+ab'),
\]
where $(a,b)$ acts by $x\mapsto ax+b$. Take
$\zeta_N=e^{2\pi i/N}$ and the positive real roots
$\theta_N=2^{1/N}$, compatibly under divisibility. An automorphism
$\sigma\in\operatorname{Gal}(\overline\Q/\Q)$ then has coordinates
\[
 \sigma(\zeta_N)=\zeta_N^a,\qquad
 \sigma(\theta_N)=\zeta_N^b\theta_N.
\]
For odd $a$, define two characters by
\[
\begin{array}{c|rrrr}
 a\bmod8&1&3&5&7\\\hline
 \chi_2(a)&1&-1&-1&1\\
 \chi_{-2}(a)&1&1&-1&-1.
\end{array}
\]
The exponent $(-1)^b$ depends on $b\bmod2$.

\subsection{The radicals of two and minus one half}
\begin{proof}[Proof of the rank-one assertions of Theorem~\ref{thm:kummer-main}]
The identity $\sqrt2=\zeta_8+\zeta_8^{-1}$ gives the necessary
relation $(-1)^b=\chi_2(a)$. To prove sufficiency, put
$L=\Q(\mu_\infty)$. We first determine precisely the radical
relations of $2$ in $L$.

If $2^j$ is an $N$th power in $L$, then $2^{j/N}\in L$, since
$L$ contains all roots of unity. Reduce $j/N$ to $u/d$ with
$\gcd(u,d)=1$. B\'ezout's identity gives
$\Q(2^{u/d})=\Q(2^{1/d})$. The polynomial $X^d-2$ is Eisenstein
at $2$. For $d>2$, its positive real root generates a real field
which omits a nonreal conjugate, so that field is not normal over
$\Q$. Every finite subfield of the abelian extension $L/\Q$ is
normal. Hence $d$ is $1$ or $2$. Conversely $1$ and $\sqrt2$
belong to $L$.

Here is the finite Kummer argument in the form needed. If a field
$F$ contains $\mu_N$, adjoining $N$th roots of elements $u_1,\ldots,u_r$
gives a Galois group embedded in $(\Z/N)^r$. Under the usual
character pairing, its annihilator consists exactly of those
$(j_1,\ldots,j_r)$ for which $\prod_i u_i^{j_i}\in F^{\times N}$:
the corresponding product of the chosen radicals is fixed exactly
in that case. The perfect pairing on the finite group then determines
the Galois image from these relations. In rank one, the preceding
paragraph shows that the translation image over $L$ is the image of
$2\Z/N$ for every $N$. Taking inverse limits gives $2\widehat\Z$.

The cyclotomic projection of the full Galois group is
$\widehat\Z^\times$. Each fibre over a multiplier is therefore a
coset of $2\widehat\Z$, and the necessary parity relation selects
that entire coset. This proves that the image is exactly
\[
 H_2=\{(a,b):(-1)^b=\chi_2(a)\}.
\]
The compatible roots of $1/2$ are $\theta_N^{-1}$, and their
translation coordinate is $-b$. Their image is the same subgroup.

At finite level $N$, if $8\mid N$ both the character and the parity
are determined by the reduced coordinates, giving index two.
If $8\nmid N$, every pair in
$(\Z/N)^\times\ltimes\Z/N$ lifts to $H_2$. For odd $N$ one can
choose the parity of $b$ freely. For $v_2(N)=1$ or $2$, the unused
residue of $a$ modulo $8$ can be chosen to give either value of
$\chi_2$; the Chinese remainder theorem imposes the remaining odd
congruences. Thus, putting $d_N=2$ when $8\mid N$ and $d_N=1$
otherwise, the finite image has order $N\varphi(N)/d_N$.

For $-1/2$, choose $r_N=\zeta_{2N}\theta_N^{-1}$.
These roots are compatible and satisfy $r_N^N=-1/2$.
The action has translation coordinate
\[
 c=t(a)-b,\qquad t(a)=(a-1)/2\in\widehat\Z.
\]
Division is well defined: $a-1$ belongs to $2\widehat\Z$ and
multiplication by $2$ is injective on $\widehat\Z$.
The identity $t(aa')=t(a)+at(a')$ shows that
$(a,b)\mapsto(a,t(a)-b)$ is an automorphism of the affine group.
Since $(-1)^{t(a)}\chi_2(a)=\chi_{-2}(a)$, its image is
\[
 H_{-2}=\{(a,c):(-1)^c=\chi_{-2}(a)\}.
\]
The same lifting argument gives the same finite orders and indices.
\end{proof}

\subsection{The three Fermat fibres}
Consider the Fermat curve
\[
 u_0^N+u_1^N+u_\infty^N=0
\]
and its map to the projective line with coordinate
$z=-u_0^N/u_\infty^N$. Over each of $z=2,-1,1/2$, choose the
coordinate whose $N$th power is the exceptional one and divide
all projective coordinates by it. The other two coordinates then
have $N$th power $-1/2$. Consequently each fibre is a product of
two copies of the same $-1/2$ Kummer torsor.

\begin{proposition}\label{prop:fermat-fibre}
With equal origins $r_N$ in both factors, the Galois image acts by
\[
 (x,y)\longmapsto(ax+c,ay+c),\qquad (a,c)\in H_{-2,N}.
\]
Its index in the scalar affine group
$(\Z/N)^2\rtimes(\Z/N)^\times$ is $Nd_N$, and its index in
$(\Z/N)^2\rtimes\operatorname{GL}_2(\Z/N)$ is
\[
 d_N N^4\prod_{\ell\mid N}(1-\ell^{-2}).
\]
Its adelic index in either ambient group is infinite. If the first
origin is instead $\zeta_N^{-1}r_N$, the first translation coordinate
is $c-a+1$ and the second remains $c$.
\end{proposition}
\begin{proof}
Equal origins give identical translation cocycles, proving the
diagonal formula. The representation on a fibre is faithful on
the rank-one image, whose order is $N\varphi(N)/d_N$.
The scalar ambient group has order $N^2\varphi(N)$, giving the
first index. For a prime power $\ell^k$, reduction from
$\operatorname{GL}_2(\Z/\ell^k)$ to
$\operatorname{GL}_2(\F_\ell)$ has kernel of order
$\ell^{4(k-1)}$, while
$|\operatorname{GL}_2(\F_\ell)|=(\ell^2-1)(\ell^2-\ell)$.
Multiplying over primes yields
\[
 |\operatorname{GL}_2(\Z/N)|
 =N^4\prod_{\ell\mid N}(1-\ell^{-1})(1-\ell^{-2}).
\]
Dividing the full affine order by the image order proves the second
formula. The unbounded indices $Nd_N$ force infinite adelic index.
For the changed origin, $x$ is replaced by $x+1$; writing $x'=x+1$, the
affine action becomes $x'\mapsto a(x'-1)+c+1$, hence the
stated coordinate.
\end{proof}
At levels $N=1,2,4,8$, the image orders are respectively $1,2,8,16$
and the scalar fibre indices are $1,2,4,16$. The first departure from
the unrestricted rank-one affine group is therefore visible at level
$8$, whereas the diagonal restriction is already visible at level $2$.

\subsection{A uniform bound for independent radicals}
The general openness assertion is classical Kummer theory. Schinzel's
abelian-binomial theorem supplies the bounded radical obstruction;
Perucca and Sgobba state it explicitly in
\cite[Theorem~15, p.~7]{PeruccaSgobba} and prove uniform degree bounds
in their Theorems~1 and~13, pp.~1 and~6--8. More recently, Perucca
\cite[Theorem~26, pp.~13--14]{Perucca2026} describes the degree-to-index
ratio for all radical levels of a finitely generated subgroup over a
number field from finitely many levels and explicit cyclotomic factors.
That result concerns the growth across all levels. We give the elementary
cohomological argument needed for the explicit saturation bound here.

Let $K$ be a number field, and let $\xi_1,\ldots,\xi_r\in K^\times$
be independent modulo roots of unity. Choose a finite set of finite
places $S$ outside which they are units. The group
\[
 U_S=\mathcal O_{K,S}^{\times}/\mu_K
\]
is a finite free abelian group by the $S$-unit theorem
\cite[Theorem~5.11, p.~90]{MilneANT}. Let $M\subset U_S$ be the
rank-$r$ subgroup spanned by their classes, and put
\[
 M^{\mathrm{sat}}=(M\otimes\Q)\cap U_S.
\]
Choose a positive integer $e$ annihilating the finite quotient
$M^{\mathrm{sat}}/M$, and put $w=|\mu_K|$ and $D=ew$.

\begin{proposition}\label{prop:general-kummer-bound}
Over $L=K(\mu_\infty)$, the inverse-limit Kummer translation image
contains $D\widehat\Z^r$. Over $K$, the scalar affine image has
index at most
\[
 [\widehat\Z^\times:\chi(G_K)]D^r
 =[K\cap\Q(\mu_\infty):\Q]D^r.
\]
It is therefore open in
$\widehat\Z^r\rtimes\widehat\Z^\times$.
For $r\ge2$ it has infinite index in
$\widehat\Z^r\rtimes\operatorname{GL}_r(\widehat\Z)$.
\end{proposition}
\begin{proof}
Fix $N$ and write $G=\operatorname{Gal}(L/K)$, an abelian group.
For a multiplicative cocycle $z:G\to\mu_N$, commutativity gives
\[
 z(t)^{\chi(s)-1}=z(s)^{\chi(t)-1}.
\]
For fixed $s$, the right side is the coboundary of $z(s)$.
Hence $\chi(s)-1$ annihilates $H^1(G,\mu_N)$ for every $s$.
The ideal of $\Z/N$ generated by these differences has annihilator
exactly $\mu_N(K)$ inside the cyclic module $\mu_N$.
Its generator is therefore $d=|\mu_N(K)|=\gcd(N,w)$.
It follows that $d$, and thus $w$, annihilates $H^1(G,\mu_N)$.
All cocycles in question are continuous; the calculation takes place
in their finite image and needs no finiteness assumption on $G$.

Suppose $u=\prod_j\xi_j^{a_j}=\alpha^N$ with $\alpha\in L$.
The cocycle $z(s)=s(\alpha)/\alpha$ has $w$th power a coboundary,
so for some $\zeta\in\mu_N$ the element
$y=\alpha^w/\zeta$ lies in $K$. Thus
\[
 y^N=u^w.
\]
Valuations outside $S$ show that $y$ is an $S$-unit, and in $U_S$
we have $N[y]=w\sum_j a_j[\xi_j]$. In particular
$[y]\in M^{\mathrm{sat}}$. Multiplication by $e$ puts $[y]$ in
$M$, whose basis is the independent list $[\xi_j]$. Comparing
coefficients in
\[
 N e[y]=D\sum_j a_j[\xi_j]
\]
gives $N\mid Da_j$ for every $j$.

By the finite Kummer pairing proved above, the annihilator of the
translation image in $(\Z/N)^r$ is contained in the kernel of
multiplication by $D$. Taking annihilators reverses this inclusion,
so the image contains $D(\Z/N)^r$. A closed subgroup of
$\widehat\Z^r$ is the intersection of the inverse images of its
finite quotient images. Consequently the inverse-limit image
contains $D\widehat\Z^r$.

The projection to multipliers is exactly $\chi(G_K)$, of index
$[K\cap\Q(\mu_\infty):\Q]$ in $\widehat\Z^\times$.
The translation subgroup has index at most $D^r$. Multiplying the
indices proves the asserted bound. When $r\ge2$, all linear parts
are scalar matrices. Modulo $\ell^k$, scalar matrices have order
$\varphi(\ell^k)$ while
$|\operatorname{GL}_r(\Z/\ell^k)|$
grows as $\ell^{r^2k}$ up to a factor independent of $k$.
Their index is unbounded, proving the last assertion. For $r=0$
the statement reduces to the cyclotomic image alone.
\end{proof}


\section{Scope and further questions}\label{sec:scope}
The all-ring kernel criterion is an exact existence condition, not a finite terminating algorithm for arbitrary infinite ring presentations. The reconstruction theorem concerns the full natural action on finite tests and does not identify a cohomology ring with that information. In the dessin classification, fixing white points removes the dilation group of an orange but leaves positive dilations for the one-loop sphere; fixing face centres removes both continuous families. The sextic cubic is the classical one credited above. The noncompact assertion requires ordinary disk faces with finite attaching cycles; it makes no claim for arbitrary asymptotic tracts. The general Kummer bound requires multiplicative independence modulo roots of unity; repeated radicals instead give the diagonal restrictions calculated in Section~\ref{sec:kummer}.

\section*{Method and records}
The mathematics and text were produced by GPT-6 Astra (ChatGPT Pro mode).
Separate instances of the same model reconstructed the main coherence,
walking-complex, cochain and compact dessin results from statements
before reading their proofs; effort settings were not exposed.
Additional statement-derived and targeted checks, with their exact
scopes and computational evidence, are listed in
\href{\recordurl}{the records}. No human has checked the work.
The text is dedicated to the public domain under CC0~1.0.
The programs \emph{Divided-power certificates},
\emph{Scalar forward-difference audit}, \emph{Sextic exact audit}
and \emph{Kummer finite-image audit} accompany the proposed
\href{\workbenchurl}{workbench}.
\begin{thebibliography}{99}
\small
\bibitem{Groth} A. Grothendieck, Montpellier manuscripts, cotes~137--138. \url{https://grothendieck.umontpellier.fr/137.pdf}; \url{https://grothendieck.umontpellier.fr/138.pdf}.
\bibitem{Buehler} T. B\"uhler, \emph{Exact Categories}, ESI preprint 2063 (14 November 2008), Theorem~A.1 p.~36 and Proposition~7.5 p.~17. \url{https://www.esi.ac.at/preprints/esi2063.pdf}.
\bibitem{Roby} N. Roby, \emph{Lois polyn\^omes et lois formelles en th\'eorie des modules}, Ann. Sci. \'Ecole Norm. Sup. (3) 80 (1963), 213--348. Theorem~IV.1, p.~266 (PDF page~55). \url{https://www.numdam.org/item/ASENS_1963_3_80_3_213_0.pdf}.
\bibitem{Krause} H. Krause, \emph{Koszul, Ringel and Serre duality for strict polynomial functors}, Compos. Math. 149 (2013). Author version, Section~2, pp.~3--5. \url{https://arxiv.org/abs/1203.0311}.
\bibitem{Horel} G. Horel, \emph{Binomial rings and homotopy theory}, arXiv:2211.02349v2 (23 March 2025), Theorem~3.7 p.~7, Section~4 p.~9 and Theorem~5.4 p.~13. \url{https://arxiv.org/abs/2211.02349v2}.
\bibitem{Hua} M. Hua, \emph{Grothendieck Montpellier. \'Enonc\'es des lectures modernis\'ees d\'emontr\'es}, version~2 (5 November 2026), pp.~165--168. \url{https://hexagonmath.org/2611.00125v2}.
\bibitem{Riehl} E. Riehl, \emph{Category Theory in Context}, author online version, Definition~5.2.10, p.~194. \url{https://emilyriehl.github.io/files/context.pdf}.
\bibitem{CLF} A. Chambert-Loir and M.~I. de Frutos-Fern\'andez, \emph{Formalizing Polynomial Laws and the Universal Divided Power Algebra}, CPP~2026, pp.~75--87. \url{https://doi.org/10.1145/3779031.3779108}.
\bibitem{Schneps} L. Schneps, \emph{Dessins d'enfants on the Riemann sphere}, in The Grothendieck Theory of Dessins d'Enfants, LMS Lecture Note Series 200 (1994), 47--78. \url{https://webusers.imj-prg.fr/~leila.schneps/Fschneps.pdf}.
\bibitem{DoldPuppe} A. Dold and D. Puppe, \emph{Homologie nicht-additiver Funktoren. Anwendungen}, Ann. Inst. Fourier \textbf{11} (1961), 201--312. Equation~(3.3), p.~219 (PDF p.~20); Theorem~3.6 and Corollary~3.8, p.~220 (PDF p.~21). \url{https://www.numdam.org/item/AIF_1961__11__201_0.pdf}.
\bibitem{CFH} A. Conway, S. Friedl and G. Herrmann, \emph{Linking forms revisited}, author version of 26 February 2018. Quotient convention p.~1; Proposition~1.2, p.~2. \url{https://friedl.app.uni-regensburg.de/papers/linking-forms-february-26-2018.pdf}.
\bibitem{StacksRegular} The Stacks Project Authors, \emph{The Stacks Project}, Lemma~15.76.14, Tag~066Z, regular noetherian rings and perfect complexes. Accessed 23 November 2026. \url{https://stacks.math.columbia.edu/tag/066Z}.
\bibitem{StacksFibres} The Stacks Project Authors, \emph{The Stacks Project}, Lemma~15.79.2, Tag~068V, recognizing perfect complexes from residue-field fibres. Accessed 23 November 2026. \url{https://stacks.math.columbia.edu/tag/068V}.
\bibitem{Mandell} M. A. Mandell, \emph{Cochains and homotopy type}, Publ. Math. IHES \textbf{103} (2006), 213--246. Main Theorem and Theorems~0.1--0.2, pp.~213--216. \url{https://www.numdam.org/article/PMIHES_2006__103__213_0.pdf}.
\bibitem{Wendl} C. Wendl, \emph{Lectures on Holomorphic Curves in Symplectic and Contact Geometry}, version~3.3 (7 May 2015), Theorem~4.2.1 and quotient classification, printed pp.~152--153. \url{https://www2.mathematik.hu-berlin.de/~wendl/pub/jhol_bookv33_2.pdf}.
\bibitem{Canas} J. Cañas, R. A. Hidalgo, F. J. Turiel and A. Viruel, \emph{Groups as automorphisms of dessins d'enfants}, Rev. R. Acad. Cienc. Exactas Fís. Nat. Ser. A Mat. \textbf{116} (2022), article~160. Theorems~1--2, article p.~3; definitions, p.~4. \url{https://doi.org/10.1007/s13398-022-01285-7}.
\bibitem{MilneANT} J. S. Milne, \emph{Algebraic Number Theory}, version~3.08 (19 July 2020), Theorem~5.11, printed p.~90. \url{https://www.jmilne.org/math/CourseNotes/ANT.pdf}.
\bibitem{PeruccaSgobba} A. Perucca and P. Sgobba, \emph{Kummer theory for number fields and the reductions of algebraic numbers}, Int. J. Number Theory \textbf{15} (2019), 1617--1633. Author version, Theorems~1, 13 and~15, pp.~1 and~6--8. \url{https://orbilu.uni.lu/bitstream/10993/35969/1/PeruccaSgobba-update.pdf}.
\bibitem{Schinzel} A. Schinzel, \emph{Abelian binomials, power residues and exponential congruences}, Acta Arith. \textbf{32} (1977), 245--274. Theorem~2, p.~245; its radical form is stated in \cite[Theorem~15, p.~7]{PeruccaSgobba}. \url{https://doi.org/10.4064/aa-32-3-245-274}.
\bibitem{Toen} B. To\"en, \emph{Le probl\`eme de la sch\'ematisation de Grothendieck revisit\'e}, Épijournal de Géométrie Algébrique \textbf{4} (2020), article~14. Theorem~4.1 p.~11; Corollaries~5.1, 5.4--5.5, pp.~14 and~16--17. \url{https://epiga.episciences.org/6822/pdf}.
\bibitem{Perucca2026} A. Perucca, \emph{Linear relations among radicals}, arXiv:2511.02498v2 (9 November 2026), Theorem~26, pp.~13--14. \url{https://arxiv.org/abs/2511.02498v2}.
\bibitem{FantechiRicolfi} B. Fantechi and A. T. Ricolfi, \emph{On the stack of 0-dimensional coherent sheaves: Motivic aspects}, Bull. London Math. Soc. \textbf{57} (2025), 1607--1649. Author version, Appendix~B, pp.~27--30. \url{https://arxiv.org/abs/2403.07859}.
\bibitem{MulasePenkava} M. Mulase and M. Penkava,
\emph{Ribbon graphs, quadratic differentials on Riemann surfaces,
and algebraic curves defined over $\overline{\Q}$}, Asian J. Math.
\textbf{2} (1998), 875--920. Definition~1.9 and classification,
pp.~880--881 (author version pp.~7--8).
\url{https://doi.org/10.4310/AJM.1998.v2.n4.a11}.
\bibitem{Richter} B. Richter, \emph{Divided power structures and chain
complexes}, Contemp. Math. \textbf{504} (2009), 237--254.
Theorems~7.1 and~7.5, pp.~248--250.
\url{https://doi.org/10.1090/conm/504/09882}.
\bibitem{Ikonicoff} S. Ikonicoff, \emph{Divided power algebras over an
operad}, Glasgow Math. J. \textbf{62} (2020), 477--517.
Author version arXiv:1712.03694v2, Definition~3.1.4 and
Theorem~3.1.8, pp.~12--13.
\url{https://arxiv.org/abs/1712.03694v2}.
\end{thebibliography}
\end{document}
